\documentclass[11pt,letterpaper]{article}
\usepackage[margin=1in]{geometry}
\usepackage[T1]{fontenc}
\usepackage{lmodern,amsmath,amssymb,amsthm,mathtools}
\usepackage{booktabs,longtable,array,enumitem,xcolor,microtype,xurl,fancyhdr}
\usepackage[colorlinks=true,linkcolor=blue!45!black,citecolor=blue!45!black,urlcolor=blue!50!black]{hyperref}
\usepackage{bookmark}
\usepackage{needspace}
\emergencystretch=3em
\setlength{\headheight}{14pt}
\pagestyle{fancy}\fancyhf{}
\fancyhead[L]{\small Resonant jets and exact identities}
\fancyhead[R]{\small ProveIt research continuation}
\fancyfoot[C]{\thepage}
\setcounter{tocdepth}{2}
\numberwithin{equation}{section}
\theoremstyle{plain}
\newtheorem{theorem}{Theorem}[section]
\newtheorem{proposition}[theorem]{Proposition}
\newtheorem{lemma}[theorem]{Lemma}
\newtheorem{corollary}[theorem]{Corollary}
\theoremstyle{definition}
\newtheorem{definition}[theorem]{Definition}
\theoremstyle{remark}
\newtheorem{remark}[theorem]{Remark}
\newtheorem{example}[theorem]{Example}
\DeclareMathOperator{\Li}{Li}
\DeclareMathOperator{\FP}{FP}
\DeclareMathOperator{\Log}{Log}
\DeclareMathOperator{\Rea}{Re}
\DeclareMathOperator{\artanh}{artanh}
\DeclareMathOperator{\Sym}{Sym}
\newcommand{\N}{\mathbb N}
\newcommand{\R}{\mathbb R}
\newcommand{\C}{\mathbb C}
\newcommand{\Q}{\mathbb Q}
\newcommand{\Z}{\mathbb Z}
\newcommand{\eps}{\varepsilon}
\newcommand{\dd}{\,\mathrm d}
\newcommand{\ii}{\mathrm i}
\newcommand{\cG}{\mathcal G}
\newcommand{\cZ}{\mathcal Z}
\newcommand{\cH}{\mathcal H}
\newcommand{\cR}{\mathcal R}
\newcommand{\laur}[1]{[\eps^{#1}]}
\newcommand{\src}[1]{\nolinkurl{#1}}
\newcommand{\poch}[2]{(#1)_{#2}}
\title{Resonant Gamma Jets and Directional Zeta Identities\\[7pt]
\large Discrete Stieltjes products, shifted polylogarithms,\\
and local formulas for higher spectral derivatives}
\author{ProveIt research continuation\\\small Prepared for Vladimir Reshetnikov}
\date{10 October 2026}
\hypersetup{pdftitle={Resonant Gamma Jets and Directional Zeta Identities},pdfauthor={ProveIt research continuation},pdfsubject={Polylogarithms, Gamma, Hurwitz zeta, generalized Stieltjes constants}}
\begin{document}
\maketitle
\begin{abstract}
We develop six compatible exact calculi beyond the current polylogarithm
manuscript. A discrete product identity evaluates convergent sums containing
arbitrarily many generalized Stieltjes functions, with an explicit finite
remainder. A holomorphic remainder for the double Hurwitz zeta function
separates its two intersecting polar divisors, yielding the complete
dependence of its constant term on a transverse direction or a curved
approach. Symmetrized nested sums admit an all-depth, all-index partition
formula. The higher spectral derivatives of a renormalized reciprocal-Gamma
product at every moving Gamma zero are then reduced to a finite Bell
algorithm and a finite triangular family of ordinary antiderivatives.
An independent shifted-gap identity gives complete reductions of
complementary-shift double polylogarithms throughout the open argument
disk. Finally, higher derivatives of polynomially weighted spectral zeta
functions satisfy finite-difference identities with explicitly local right
sides; a six-term second-derivative identity has a Vandermonde evaluation.
It belongs to an alternating Vandermonde family with any number of factors,
including exact identities for fifth and higher spectral derivatives.
All stated identities are proved, with domains and normalizations. The
classical symmetric-sum, parity, and multiplicative-anomaly mechanisms are
credited explicitly. Numerical checks are reproducible diagnostics, not
substitutes for proofs. The short Gaussian $S_6$ and $S_8$ evaluations and
arithmetic minimality questions remain open.
\end{abstract}
\vspace{5pt}
\noindent\textbf{Source baseline.} Repository revision
\src{16c7e342d7a4a15f59911f322eb90b7f96c2a8b5}; the canonical manuscript and
all five archives present in \src{docs/incoming} at that revision.
\clearpage
\tableofcontents
\clearpage
% BEGIN sections/01-scope.tex
\section{Scope, provenance, and conventions}\label{sec:scope}

The immediate source is the collective manuscript \emph{Polylogarithms and
their Arithmetic Bridges}, together with the five incoming archives in the
pinned revision \cite{repo,nested,gauss,mellin,stharm,translation}. The
manuscript already contains a substantial theory of periodic Stieltjes
correlations, distributional derivatives, dilation contacts, and
log-Gamma convolutions. Repeating those results would not answer the
request for a research continuation. Here the focus is on discrete
translation, independent spectral directions, moving zeros, and
nonperiodic shifted sums.

The incoming \emph{Nested Harmonic Jets} report asks, among other things,
for non-diagonal spectral directions, path-dependent multiple Laurent
finite parts, and an explicit algorithm for higher spectral values at
moving Gamma zeros. We answer the algorithmic part of the last question,
give a complete transverse-path calculation at depth two, and give every
Laurent coefficient of the fully symmetrized family at arbitrary depth.
The unsymmetrized higher-depth problem and minimality of the resulting
constant classes are not settled by these results.

\subsection{What is established here}

\begin{longtable}{@{}>{\raggedright\arraybackslash}p{.24\textwidth}>{\raggedright\arraybackslash}p{.49\textwidth}>{\raggedright\arraybackslash}p{.21\textwidth}@{}}
\toprule Topic & Exact conclusion & Location\\\midrule
\endhead
Discrete Stieltjes products & One explicit correction term evaluates the
grouped sum for every number of factors and all Laurent indices; the
finite-cutoff remainder is exact. & Section~\ref{sec:discrete}\\[4pt]
Intersecting nested poles & A holomorphic two-variable remainder, all its
Taylor coefficients, straight-direction finite parts, and the cubic-jet
correction for curved paths. & Section~\ref{sec:direction}\\[4pt]
Symmetrized mixed jets & A finite partition formula for every Laurent
coefficient in every depth, including an explicit depth-three constant
term. & Section~\ref{sec:symmetric}\\[4pt]
Moving Gamma zeros & A Bell-polynomial formula for every derivative, a
closed second derivative, and a finite antiderivative construction for
all necessary coordinates. & Section~\ref{sec:roots}\\[4pt]
Shifted double polylogarithms & A disk identity in two independent shifts;
complete reductions at complementary shifts, including rational-shift
specializations. & Section~\ref{sec:gap}\\[4pt]
Higher spectral products & Polynomiality in normalized exponents, local
finite differences, alternating Vandermonde evaluations for any number
of factors, and logarithmic multiplicities. & Section~\ref{sec:spectral}\\
\bottomrule
\end{longtable}

These are proved extensions developed for this continuation. This wording
does not assert historical priority for every displayed specialization.
The symmetric-sum theorem is classical \cite{hoffman}; Laurent expansions
of multiple zeta functions have an established literature
\cite{mow,guozhang}; broad shifted cyclotomic symmetry is treated in
\cite{rui,xu}; and first-derivative multiplicative anomalies are classical
in substance \cite{dowker,gurel,mizuno}. The contributions should be
assessed as explicit formulas, proofs, normalizations, and extensions of
the repository's available results, with those antecedents retained.

\subsection{Basic notation}

For $a>0$, let
\begin{equation}\label{eq:stieltjes}
 \zeta(1+u,a)=\frac1u+
 \sum_{j=0}^{\infty}\frac{(-1)^j}{j!}\gamma_j(a)u^j,
 \qquad g_a(u)=\zeta(1+u,a)-\frac1u.
\end{equation}
The function $g_a$ is entire in $u$. Thus $\gamma_0(a)=-\psi(a)$,
where $\psi=\Gamma'/\Gamma$, and $\gamma_j=\gamma_j(1)$.
Spectral derivatives of Hurwitz zeta are written $\zeta^{(r)}(s,a)$;
derivatives of $\gamma_j(a)$ refer to $a$. The standard shift and
differentiation identities are \cite[\S25.11]{dlmf}
\begin{align}
 \gamma_j(a+1)&=\gamma_j(a)-\frac{\log^j a}{a},\label{eq:gshift}\\
 \partial_a\zeta(s,a)&=-s\zeta(s+1,a),\label{eq:zshiftder}\\
 \psi^{(k)}(a)&=(-1)^{k+1}k!\zeta(k+1,a)\quad(k\ge1).
\end{align}
All logarithms of positive real numbers are real. When a complex argument
is introduced, its branch is specified locally. No unspecified principal
logarithm of a Gamma quotient is used.

For a meromorphic one-variable germ $f(\eps)$,
\begin{equation}
 \FP_{\eps=0}f=\laur0f.
\end{equation}
The coefficient depends on the parameter $\eps$. A function of several
variables need not have a path-independent finite part. Nor is the
constant Laurent coefficient automatically an algebra homomorphism.
We retain every regulator until after multiplication or composition.

The exponential Bell polynomials $Y_r$ are defined by
\begin{equation}\label{eq:bell}
 \exp\left(\sum_{j\ge1}x_j\frac{t^j}{j!}\right)
 =\sum_{r\ge0}Y_r(x_1,\ldots,x_r)\frac{t^r}{r!}.
\end{equation}
In particular $Y_0=1$, $Y_1=x_1$, $Y_2=x_1^2+x_2$, and
$Y_3=x_1^3+3x_1x_2+x_3$. Generalized harmonic numbers are
$H_N^{(j)}=\sum_{k=1}^N k^{-j}$, with $H_0^{(j)}=0$.

\subsection{What the verification does and does not establish}

The proofs below give the analytic statements. The accompanying programs
check independently constructed finite coefficients, special-function
values, quadrature formulas, and the signs and scale factors in the
spectral identities. Exact symbolic assertions are separated from
floating-point diagnostics. The package makes no claim of proof-assistant
formalization or interval certification. In particular, a small residual
does not establish the short $S_6$ or $S_8$ evaluations.

% END sections/01-scope.tex
% BEGIN sections/02-discrete-products.tex
\section{An all-degree discrete Stieltjes product identity}\label{sec:discrete}

The elementary shift law \eqref{eq:gshift} contains a useful exact
nonlinear summation calculus. The obstacle is that a product of Stieltjes
functions grows as a power of $\log x$. A single additional Stieltjes
function cancels both its leading term and its first inverse-power term.
This makes a telescoping identity into an ordinary convergent series.

\begin{lemma}[The two terms needed for product cancellation]\label{lem:gasym}
For each fixed integer $p\ge0$, as $x\to+\infty$,
\begin{equation}\label{eq:gasym}
 \gamma_p(x)=-\frac{\log^{p+1}x}{p+1}
       +\frac{\log^p x}{2x}
       +O\left(\frac{1+\log^p x}{x^2}\right).
\end{equation}
The expansion can be differentiated any fixed number of times in $x$,
with the corresponding differentiated error estimate.
\end{lemma}
\begin{proof}
Use Euler--Maclaurin for $\zeta(1+u,x)$, uniformly on a small circle
about $u=0$:
\[
 \zeta(1+u,x)=\frac{x^{-u}}u+
 \frac{x^{-1-u}}2+\frac{1+u}{12}x^{-2-u}
 +O(x^{-4-\Rea u}).
\]
The Euler--Maclaurin remainder is holomorphic in $u$ there after the
displayed pole is removed. Its integral form permits any fixed number
of $u$ and $x$ derivatives: each $u$ derivative inserts a power of a
logarithm, and each $x$ derivative increases the inverse-power decay.
Coefficient extraction therefore gives \eqref{eq:gasym} with its
displayed logarithmic error, and the same argument gives every fixed
differentiated error estimate. This is a fixed-index use of the standard
Hurwitz expansion, not an assertion uniform in growing $p$.
\end{proof}

Fix $r\ge2$ and integers $p_1,\ldots,p_r\ge0$. Put
\begin{equation}\label{eq:product-K}
 K=\sum_{i=1}^r(p_i+1),\qquad
 A=\prod_{i=1}^r(p_i+1),\qquad c=\frac{(-1)^rK}{A},
\end{equation}
and define
\begin{equation}\label{eq:product-F}
 F_{\boldsymbol p}(x)=\prod_{i=1}^r\gamma_{p_i}(x)
                      +c\gamma_{K-1}(x).
\end{equation}
The finite subset formula below is practical even when the indices
$p_i$ are unequal.

\begin{theorem}[Arbitrary products, one correction, exact remainder]
\label{thm:discrete}
For $a>0$, set $x_n=n+a$ and
\begin{align}\label{eq:product-summand}
 W_{\boldsymbol p}(x)
 ={}&\sum_{\varnothing\ne I\subseteq\{1,\ldots,r\}}
 (-1)^{|I|+1}\frac{\log^{\sum_{i\in I}p_i}x}{x^{|I|}}
             \prod_{j\notin I}\gamma_{p_j}(x)
       +c\frac{\log^{K-1}x}{x}.
\end{align}
Then the grouped summand is absolutely summable, and
\begin{equation}\label{eq:allproduct}
 \boxed{\displaystyle
 \sum_{n=0}^{\infty}W_{\boldsymbol p}(n+a)
 =\prod_{i=1}^r\gamma_{p_i}(a)+c\gamma_{K-1}(a).}
\end{equation}
For every integer $N\ge0$ one has the exact finite identity
\begin{equation}\label{eq:product-tail}
 \sum_{n=0}^{N-1}W_{\boldsymbol p}(n+a)
 =F_{\boldsymbol p}(a)-F_{\boldsymbol p}(N+a).
\end{equation}
In particular the remaining infinite tail is precisely
$F_{\boldsymbol p}(N+a)$.
\end{theorem}
\begin{proof}
Write $h_i(x)=\log^{p_i}x/x$. The product expansion and
\eqref{eq:gshift} give
\[
 \prod_i\gamma_{p_i}(x)-\prod_i\gamma_{p_i}(x+1)
 =\sum_{\varnothing\ne I}(-1)^{|I|+1}
         \prod_{i\in I}h_i(x)\prod_{j\notin I}\gamma_{p_j}(x).
\]
Adding $c$ times the shift law at index $K-1$ proves
$W_{\boldsymbol p}(x)=F_{\boldsymbol p}(x)-F_{\boldsymbol p}(x+1)$,
and therefore \eqref{eq:product-tail}.

For completeness, the cancellation at infinity is stronger than merely
removing a divergent logarithm. Lemma~\ref{lem:gasym} gives
\[
 \prod_i\gamma_{p_i}(x)
 =\frac{(-1)^r}{A}\log^Kx
  +\frac{(-1)^{r-1}K}{2A}\frac{\log^{K-1}x}{x}
  +O\left(\frac{1+\log^{K-1}x}{x^2}\right).
\]
Both displayed terms cancel against $c\gamma_{K-1}(x)$. Thus
$F_{\boldsymbol p}(x)=O((1+\log^{K-1}x)/x^2)$, and its derivative is
$O((1+\log^{K-1}x)/x^3)$. Integrating the derivative from $x$ to $x+1$
proves absolute summability of the grouped $W_{\boldsymbol p}$. Passing
to the limit in the exact finite identity proves \eqref{eq:allproduct}.
\end{proof}

\begin{remark}[Grouping is part of the statement]
The individual terms in \eqref{eq:product-summand} usually have divergent
series. Equation~\eqref{eq:allproduct} is an ordinary convergent sum of
the displayed grouped expression. It does not license replacing each
separate divergent series by an independently chosen finite part.
\end{remark}

\subsection{A symmetric two-index summation formula}

For $m,n\ge0$ define the absolutely convergent series
\begin{equation}\label{eq:Sdef}
 S_{m,n}(a)=\sum_{k=0}^{\infty}
 \frac{\log^m(k+a)}{k+a}
 \left\{\gamma_n(k+a)+\frac{\log^{n+1}(k+a)}{n+1}
                  -\frac{\log^n(k+a)}{2(k+a)}\right\}.
\end{equation}
The braces are $O((1+\log^n(k+a))/(k+a)^2)$.

\begin{corollary}[All-index Stieltjes summation]\label{cor:Stwo}
For $a>0$ and $m,n\ge0$,
\begin{equation}\label{eq:Stwo}
 \boxed{\displaystyle S_{m,n}(a)+S_{n,m}(a)
 =\gamma_m(a)\gamma_n(a)+
 \left(\frac1{m+1}+\frac1{n+1}\right)\gamma_{m+n+1}(a).}
\end{equation}
In particular,
\begin{equation}\label{eq:Ssquare}
 S_{n,n}(a)=\frac{\gamma_n(a)^2}{2}
                      +\frac{\gamma_{2n+1}(a)}{n+1}.
\end{equation}
\end{corollary}
\begin{proof}
Apply Theorem~\ref{thm:discrete} with $r=2$, $p_1=m$, and $p_2=n$.
The two copies of the last term in the braces of \eqref{eq:Sdef}
combine into the equality-block correction
$-\log^{m+n}x/x^2$ in \eqref{eq:product-summand}.
\end{proof}

The first cases already connect the three types of functions sought in
the research programme:
\begin{align}
 \sum_{k=0}^{\infty}
 \frac{\log(k+a)-\psi(k+a)-1/(2(k+a))}{k+a}
 &=\frac{\psi(a)^2}{2}+\gamma_1(a),\label{eq:digamma-square}\\
 \sum_{k=0}^{\infty}\frac{\log(k+a)}{k+a}
 \left\{\gamma_1(k+a)+\frac{\log^2(k+a)}2
           -\frac{\log(k+a)}{2(k+a)}\right\}
 &=\frac{\gamma_1(a)^2+\gamma_3(a)}2.\label{eq:gamma1-square}
\end{align}
The first is a useful normalization check; the second has exactly the
same proof, with no numerical relation search.

\subsection{Cubic and higher polygamma identities}

Taking every $p_i=0$ gives $K=r$, $A=1$, and
\begin{equation}\label{eq:all-digamma}
 \sum_{k=0}^{\infty}
 \left\{(-\psi(x_k))^r-\left(-\psi(x_k)-\frac1{x_k}\right)^r
                  +(-1)^rr\frac{\log^{r-1}x_k}{x_k}\right\}
 =(-\psi(a))^r+(-1)^rr\gamma_{r-1}(a).
\end{equation}
For example,
\begin{equation}\label{eq:cubic-digamma}
 \sum_{k=0}^{\infty}\left\{
 \frac{3\psi(x_k)^2-3\log^2x_k}{x_k}
 +\frac{3\psi(x_k)}{x_k^2}+\frac1{x_k^3}\right\}
 =-\psi(a)^3-3\gamma_2(a).
\end{equation}
At an integer $a=N$, the recurrence
$\psi(N+k)=H_{N+k-1}-\gamma$ turns this into an exact harmonic-number
sum; at rational $a$ it gives the corresponding shifted harmonic sum.

Every identity in this section can be differentiated with respect to $a$
on compact subintervals of $(0,\infty)$. Indeed the differentiated
Euler--Maclaurin estimate gives locally uniform summability of all fixed
derivative orders. For reference,
\begin{equation}\label{eq:gamma-derivative}
 \gamma_n'(a)=(-1)^{n+1}
 \left\{\zeta^{(n)}(2,a)+n\zeta^{(n-1)}(2,a)\right\},
\end{equation}
where the second term is omitted for $n=0$. This follows by matching
coefficients in \eqref{eq:zshiftder}. Thus the differentiated identities
remain in an explicitly stated zeta--polygamma vocabulary.

There is also an exact primitive statement. On every compact interval
$[a,b]\subset(0,\infty)$,
\begin{equation}
 \sum_{n=0}^{\infty}\int_a^b W_{\boldsymbol p}(n+t)\dd t
 =\int_a^b\left\{\prod_i\gamma_{p_i}(t)+c\gamma_{K-1}(t)\right\}\dd t.
\end{equation}
This follows by locally uniform convergence. It fixes the integration
constant through a definite interval, instead of silently assigning a
normalization to an indefinite primitive.

% END sections/02-discrete-products.tex
% BEGIN sections/03-directional-hurwitz.tex
\section{Directional and curved finite parts of double Hurwitz zeta}
\label{sec:direction}

The incoming report asks what changes when the exponent perturbations
are independent, and whether its diagonal finite parts depend on the
approach path. At depth two these questions have a particularly explicit
answer. The meromorphic continuation itself belongs to the established
multiple-zeta theory \cite{mow}. The result below identifies its full
polar germ in the present normalization, then computes the resulting
path dependence and all regular Taylor coefficients.

For $a>0$ define
\begin{equation}\label{eq:doublehurwitz}
 T(s,t;a)=\sum_{n>m\ge0}(n+a)^{-s}(m+a)^{-t}.
\end{equation}
The defining series is absolutely convergent when
$\Rea s>1$ and $\Rea(s+t)>2$. Its order convention is essential: the
outer, larger index carries $s$. Put
\[
 g_a(u)=\zeta(1+u,a)-\frac1u
       =\sum_{j\ge0}g_j(a)u^j,
 \qquad g_j(a)=\frac{(-1)^j\gamma_j(a)}{j!}.
\]

\begin{theorem}[Exact polar germ at $(1,1)$]\label{thm:doublegerm}
There is a function $H_a(u,v)$ holomorphic near $(0,0)$ such that
\begin{equation}\label{eq:doublegerm}
 \boxed{T(1+u,1+v;a)
   =\frac1{u(u+v)}+\frac{g_a(v)}u+H_a(u,v).}
\end{equation}
It satisfies
\begin{align}
 H_a(u,v)+H_a(v,u)
   &=g_a(u)g_a(v)-\zeta(2+u+v,a),\label{eq:Hsym}\\
 H_a(0,0)&=C_a:=\frac{\gamma_0(a)^2-\zeta(2,a)}2.
 \label{eq:H00}
\end{align}
Thus the two local polar divisors are $u=0$ and $u+v=0$; in particular
there is no separate polar divisor $v=0$.
\end{theorem}

\begin{proof}
Write $x=n+a$, and use the inner finite sum to obtain, first in the
absolute-convergence domain,
\[
 T(1+u,1+v;a)=\zeta(1+u,a)\zeta(1+v,a)
       -\sum_{n\ge0}x^{-1-u}\zeta(1+v,x).
\]
The pole-subtracted Euler--Maclaurin remainder
\[
 E(v,x)=\zeta(1+v,x)-\frac{x^{-v}}v-\frac12x^{-1-v}
\]
is holomorphic in $v$ near zero and is
$O(x^{-2-\Rea v})$ uniformly on compact $v$-sets there. Hence
\[
 R_a(u,v)=\sum_{n\ge0}(n+a)^{-1-u}E(v,n+a)
\]
converges normally in a neighborhood of $(0,0)$, as do its fixed
derivatives. We obtain
\begin{equation}\label{eq:doubleEM}
 T=\zeta(1+u,a)\zeta(1+v,a)
   -\frac{\zeta(1+u+v,a)}v
   -\frac12\zeta(2+u+v,a)-R_a(u,v).
\end{equation}
The apparent singularity at $v=0$ cancels. Substitution of
$\zeta(1+u,a)=u^{-1}+g_a(u)$ gives \eqref{eq:doublegerm}, with
\begin{equation}\label{eq:Hrepresentation}
 H_a(u,v)=g_a(u)g_a(v)
  +\frac{g_a(u)-g_a(u+v)}v
  -\frac12\zeta(2+u+v,a)-R_a(u,v).
\end{equation}
The divided difference is holomorphic, since its numerator vanishes at
$v=0$. Finally, the finite-sum stuffle identity continues to
\[
 T(s,t;a)+T(t,s;a)=\zeta(s,a)\zeta(t,a)-\zeta(s+t,a).
\]
Insert \eqref{eq:doublegerm} to get \eqref{eq:Hsym} and then
\eqref{eq:H00}. All steps are identities of meromorphic germs, so no
choice of a one-dimensional limiting path is hidden in the proof.
\end{proof}

\subsection{Straight directions and their exact discrepancy}

Here $\FP_{\eps=0}$ means the coefficient of $\eps^0$ in the Laurent
expansion after the specified substitution. It is not a joint value at
$(u,v)=(0,0)$.

\begin{corollary}[Every transverse straight direction]\label{cor:direction}
Let $c,d\in\C$, with $c\ne0$ and $c+d\ne0$. Then
\begin{align}\label{eq:direction-expansion}
 T(1+c\eps,1+d\eps;a)
  =\frac1{c(c+d)\eps^2}
     +\frac{\gamma_0(a)}{c\eps}
     +C_a-\frac dc\gamma_1(a)+O(\eps),
\end{align}
and in particular
\begin{equation}\label{eq:direction-FP}
 \boxed{\FP_{\eps=0}T(1+c\eps,1+d\eps;a)
   =\frac{\gamma_0(a)^2-\zeta(2,a)}2-\frac dc\gamma_1(a).}
\end{equation}
The case $d=0$ is included. The difference between two admissible
directions is exactly
$-\gamma_1(a)(d/c-d'/c')$.
\end{corollary}
\begin{proof}
Substitute in \eqref{eq:doublegerm}, expand $g_a(d\eps)$ through its
linear term, and use \eqref{eq:H00}.
\end{proof}

For $c=d=1$, this agrees with the diagonal identity
$T(s,s;a)=(\zeta(s,a)^2-\zeta(2s,a))/2$.
For $d=0$, the finite part is $C_a$ instead. A curve contained in
$u=0$ or $u+v=0$ cannot be evaluated by restriction of the displayed
meromorphic function without an additional prescription. Curves tangent
to those divisors can have higher pole orders and are outside the
transversality statement.

\subsection{Curvature and reparametrization also matter}

The tangent direction alone does not determine the constant Laurent
coefficient. The double pole reads the cubic jet of a curve.

\begin{theorem}[Finite part along an analytic transverse curve]
\label{thm:curved}
Suppose
\[
 u=c_1\eps+c_2\eps^2+c_3\eps^3+O(\eps^4),\qquad
 v=d_1\eps+d_2\eps^2+d_3\eps^3+O(\eps^4),
\]
where $c_1\ne0$ and $A:=c_1+d_1\ne0$. Set
$B=c_2+d_2$ and $C=c_3+d_3$. Then
\begin{align}\label{eq:curved-FP}
 \FP_{\eps=0}T(1+u,1+v;a)
 ={}&C_a-\gamma_1(a)\frac{d_1}{c_1}
             -\gamma_0(a)\frac{c_2}{c_1^2}\notag\\
 &+\frac1{c_1A}\left\{
       \frac{c_2^2}{c_1^2}+\frac{B^2}{A^2}
       +\frac{c_2B}{c_1A}-\frac{c_3}{c_1}-\frac CA
                         \right\}.
\end{align}
\end{theorem}
\begin{proof}
The constant term of the holomorphic part is $C_a$. In $g_a(v)/u$,
the constant terms are $-\gamma_0(a)c_2/c_1^2$ and
$-\gamma_1(a)d_1/c_1$. Expanding both reciprocal factors in
$1/[u(u+v)]$ to second order after factoring out $c_1A\eps^2$
gives the braces. Terms of order four in the curve cannot contribute.
\end{proof}

For example, $u=\eps$, $v=\eps+\lambda\eps^2$ gives
\begin{equation}\label{eq:curve-example}
 \FP T(1+\eps,1+\eps+\lambda\eps^2;a)
       =C_a-\gamma_1(a)+\frac{\lambda^2}{8}.
\end{equation}
Thus even curves with the same tangent as the diagonal have different
finite parts. This is an explicit answer to the path-dependence question,
not merely an assertion that multivariate regularization can be ambiguous.
Even the same geometric diagonal, reparametrized by
$u=v=\eps+\lambda\eps^2+\mu\eps^3$, has finite part
\[
 C_a-\gamma_1(a)-\lambda\gamma_0(a)+\frac{3\lambda^2-2\mu}{2}.
\]
This also follows directly from the exact diagonal stuffle identity,
and independently checks the cubic-jet terms in \eqref{eq:curved-FP}.

\subsection{All regular Taylor coefficients in Stieltjes coordinates}

The convergent sums $S_{m,n}(a)$ from \eqref{eq:Sdef} give a direct
coordinate system for the complete holomorphic remainder.

\begin{theorem}[The complete two-variable regular jet]\label{thm:Hcoeff}
Write $H_a(u,v)=\sum_{m,n\ge0}h_{m,n}(a)u^mv^n$. For all $m,n\ge0$,
\begin{align}\label{eq:Hcoeff}
 h_{m,n}(a)
 ={}&\frac{(-1)^{m+n}}{m!n!}
 \left\{\gamma_m(a)\gamma_n(a)
          +\frac{\gamma_{m+n+1}(a)}{n+1}-S_{m,n}(a)\right\}\notag\\
 &-\frac{\zeta^{(m+n)}(2,a)}{2m!n!}.
\end{align}
For a transverse straight direction, every regular coefficient is
therefore explicit:
\begin{equation}\label{eq:direction-all}
 [\eps^q]T(1+c\eps,1+d\eps;a)
 =\frac{(-1)^{q+1}\gamma_{q+1}(a)d^{q+1}}{c(q+1)!}
       +\sum_{m+n=q}h_{m,n}(a)c^md^n,
 \qquad q\ge0.
\end{equation}
\end{theorem}
\begin{proof}
Normal convergence permits coefficient extraction in $R_a$. The
coefficient of $v^n$ in $E(v,x)$ is
\[
 \frac{(-1)^n}{n!}\left\{
 \gamma_n(x)+\frac{\log^{n+1}x}{n+1}-\frac{\log^n x}{2x}\right\}.
\]
Consequently $[u^mv^n]R_a=(-1)^{m+n}S_{m,n}/(m!n!)$.
The coefficient of the divided difference in
\eqref{eq:Hrepresentation} is
$-\binom{m+n+1}{n+1}g_{m+n+1}(a)$. Combining these two observations
with the coefficients of $g_a(u)g_a(v)$ and
$\zeta(2+u+v,a)/2$ proves \eqref{eq:Hcoeff}. Substitution in
\eqref{eq:doublegerm} proves \eqref{eq:direction-all}.
\end{proof}

Taking symmetric coefficients in \eqref{eq:Hsym} and using
\eqref{eq:Hcoeff} recovers \eqref{eq:Stwo}. Thus the nonlinear
Stieltjes summation law has two independent proofs: discrete telescoping
and the meromorphic double-zeta stuffle identity. This compatibility is
also an effective sign and normalization check.


% END sections/03-directional-hurwitz.tex
% BEGIN sections/04-symmetric-hurwitz.tex
\section{Every directional Laurent coefficient of fully symmetrized depth}
\label{sec:symmetric}

At depth greater than two, resolving one ordered multiple-zeta germ
requires additional information. Complete symmetrization nevertheless
admits an exact answer at every depth and every Laurent order. The
algebraic input is Hoffman's classical symmetric-sum principle
\cite{hoffman}; the present use records independent perturbation
directions and the resulting Stieltjes coefficients explicitly.

For $a>0$ and $d\ge1$ set
\[
 Z_d(s_1,\ldots,s_d;a)
   =\sum_{n_1>\cdots>n_d\ge0}\prod_{i=1}^d(n_i+a)^{-s_i},
\]
initially in its domain of absolute convergence. We sum over all
\emph{labelled} permutations,
\[
 \Sym Z_d(\boldsymbol s;a)
   =\sum_{\sigma\in S_d}Z_d(s_{\sigma(1)},\ldots,s_{\sigma(d)};a).
\]
Repeated exponent values still occur with their permutation
multiplicity. Let $\Pi_d$ denote the set partitions of $\{1,\ldots,d\}$.

\begin{proposition}[Classical partition identity]\label{prop:partition}
As an identity of meromorphic functions,
\begin{equation}\label{eq:partition}
 \Sym Z_d(\boldsymbol s;a)
 =\sum_{\pi\in\Pi_d}(-1)^{d-|\pi|}
      \prod_{B\in\pi}\left[(|B|-1)!\,
                \zeta\left(\sum_{i\in B}s_i,a\right)\right].
\end{equation}
\end{proposition}
\begin{proof}
First truncate all indices to $0\le n<N$. The symmetrized ordered
sum is the sum over all labelled tuples with distinct indices.
Inclusion--exclusion on equality patterns is M\"obius inversion in the
partition lattice. Its coefficient for a block $B$ is
$(-1)^{|B|-1}(|B|-1)!$. Each unrestricted equality block contributes
$\sum_{n<N}(n+a)^{-\sum_{i\in B}s_i}$. This proves the finite identity.
Letting $N\to\infty$ in the absolute-convergence region gives
\eqref{eq:partition}; meromorphic continuation gives the asserted
identity. Alternatively, the right side itself supplies the
continuation of the symmetrized function used below.
\end{proof}

\begin{theorem}[All-depth, all-order directional coefficient algorithm]
\label{thm:symcoeff}
Let $c_1,\ldots,c_d\in\C\setminus\{0\}$ and put
$s_i=1+c_i\eps$. For a block $B$, write $b=|B|$ and
$c_B=\sum_{i\in B}c_i$. Define the coefficient array
\begin{align}
 A_{\{i\}}(-1)&=\frac1{c_i},&
 A_{\{i\}}(k)&=\frac{(-1)^k\gamma_k(a)c_i^k}{k!}\quad(k\ge0),
 \label{eq:singletoncoeff}\\
 A_B(k)&=\frac{\zeta^{(k)}(b,a)c_B^k}{k!}\quad(k\ge0, b\ge2).
 \label{eq:blockcoeff}
\end{align}
All other entries are zero. Then, for every integer $q\ge-d$,
\begin{equation}\label{eq:symcoeff}
 \boxed{
 [\eps^q]\Sym Z_d(1+c_1\eps,\ldots,1+c_d\eps;a)
 =\sum_{\pi\in\Pi_d}(-1)^{d-|\pi|}
       \prod_{B\in\pi}(|B|-1)!
       \sum_{\substack{(k_B)_{B\in\pi}\\\sum_Bk_B=q}}
                     \prod_{B\in\pi}A_B(k_B).}
\end{equation}
For each fixed $d,q$ the displayed expression is finite. Vanishing
block sums $c_B$ with $|B|\ge2$ cause no singularity. The constant
coefficient uses only $\gamma_0(a),\ldots,\gamma_{d-1}(a)$ and
finitely many derivatives $\zeta^{(j)}(b,a)$ at integers $2\le b\le d$.
\end{theorem}
\begin{proof}
For a singleton, the zeta factor in \eqref{eq:partition} has exactly
the Laurent coefficients \eqref{eq:singletoncoeff}. For a block of
size at least two it is regular at $\eps=0$, with Taylor coefficients
\eqref{eq:blockcoeff}. Multiplication proves \eqref{eq:symcoeff}.
There are finitely many partitions, each index has lower bound $-1$
or $0$, and their sum is prescribed; hence only finitely many tuples
can contribute. For $q=0$, a singleton index is at most $d-1$,
which proves the Stieltjes-index assertion.
\end{proof}

This is an explicit finite algorithm: enumerate set partitions,
generate each block's coefficient list to the required degree, and
multiply the short Laurent polynomials. No nested infinite sums or
unspecified finite parts remain in the output.

\subsection{Depth two and depth three in closed form}

For depth two, the constant coefficient is
\begin{equation}\label{eq:sym2}
 \FP\Sym Z_2
 =\gamma_0(a)^2-\zeta(2,a)
        -\gamma_1(a)\left(\frac{c_1}{c_2}+\frac{c_2}{c_1}\right),
\end{equation}
in agreement with the sum of the two ordered values in
\eqref{eq:direction-FP} whenever both ordered restrictions are
transverse. Formula \eqref{eq:sym2} itself also applies when
$c_1+c_2=0$: the polar divisor of each ordered summand has canceled
in the symmetrized meromorphic germ before restriction.

For depth three let $c_\Sigma=c_1+c_2+c_3$. Direct expansion gives
\begin{align}\label{eq:sym3}
 \FP\Sym Z_3={}&\gamma_0(a)^3
   -\gamma_0(a)\gamma_1(a)\sum_{i\ne j}\frac{c_j}{c_i}
   +\frac{\gamma_2(a)}2\sum_{i=1}^3
                         \frac{c_i^2}{\prod_{j\ne i}c_j}\notag\\
 &-3\gamma_0(a)\zeta(2,a)
   -\zeta'(2,a)\sum_{i=1}^3\frac{c_\Sigma-c_i}{c_i}
   +2\zeta(3,a).
\end{align}
At $c_1=c_2=c_3=1$ this becomes
\begin{equation}\label{eq:sym3diagonal}
 \gamma_0(a)^3-6\gamma_0(a)\gamma_1(a)+\frac32\gamma_2(a)
 -3\gamma_0(a)\zeta(2,a)-6\zeta'(2,a)+2\zeta(3,a).
\end{equation}
The corresponding single diagonal depth-three finite part is one
sixth of \eqref{eq:sym3diagonal}; the factor $3!$ is required by our
labelled convention.

\subsection{Why coefficient extraction is not automatically a renormalization}

The Laurent constant of a product is generally different from the
product of the Laurent constants. For example,
\begin{equation}\label{eq:FPnonmult}
 \FP_{\eps=0}\bigl\{\zeta(1+c\eps,a)\zeta(1+d\eps,a)\bigr\}
 =\gamma_0(a)^2-\gamma_1(a)(c/d+d/c),
\end{equation}
whereas the product of the two individual constant coefficients is
$\gamma_0(a)^2$. Therefore one cannot assert that a collection of
directional finite parts preserves the stuffle algebra simply because
the meromorphic functions satisfy stuffle identities. Algebra-preserving
renormalizations require a compatible subtraction prescription; this is
the subject of, for example, Guo and Zhang \cite{guozhang}.

Theorems~\ref{thm:doublegerm} and \ref{thm:symcoeff} settle two precise
parts of the incoming questions: every ordered depth-two transverse
direction, and every fully symmetrized direction at arbitrary depth.
They do not give each individual ordered germ at depth three and higher,
nor do they identify a canonical renormalization independent of a chosen
path and subtraction convention.


% END sections/04-symmetric-hurwitz.tex
% BEGIN sections/05-gamma-roots.tex
% Research fragment for the parent article. No standalone preamble.
% Source antecedent: incoming Nested Harmonic Jets, article.tex,
% proposition prop:zeroflow and Further research question 7.
% Prefix mz: avoids conflicts with the source manuscript labels.

\section{All spectral derivatives at the moving Gamma zeros}
\label{mz:sec:zeros}\label{sec:roots}

The incoming \emph{Nested Harmonic Jets} report \cite{nested} introduces a holomorphic
spectral product, evaluates its first spectral derivative at each zero of
its Gamma slice, and asks for an explicit construction at every derivative
order. We give such a construction. Its coefficients have a finite
triangular description by normalized Stieltjes antiderivatives. This solves
the algorithmic part of that question; it gives a sufficient coordinate
system rather than an arithmetic minimality theorem for constants.

The Weierstrass Gamma product, Bell polynomials, Stirling numbers, and
Mellin transforms used below are classical. The contribution here is their
specific combination into the all-order zero evaluation, its finite
Stieltjes primitive closure, and a harmonic-number coefficient algorithm.
The elementary digamma integral appearing at second order is not claimed
as a new evaluation in isolation.

Fix $a>0$, put $x_n=n+a$, and adopt
\[
 \zeta(1+u,a)=\frac1u+
       \sum_{r\ge0}\frac{(-1)^r\gamma_r(a)}{r!}u^r,
 \qquad \gamma_r=\gamma_r(1),\qquad \gamma_0(a)=-\psi(a).
\]
Use the source product and its specified counterterm:
\begin{align}
 \mathcal G(s,z;a)
 &=e^{z(\zeta(s,a)-1/(s-1))}
       \prod_{n\ge0}(1+zx_n^{-s})e^{-zx_n^{-s}},
       \qquad \Re s>\tfrac12,\label{mz:G}\\
 F_R(z;a)&=\left.\partial_s^R\mathcal G(s,z;a)\right|_{s=1}.
 \label{mz:F}
\end{align}
The pole-subtracted expression in \eqref{mz:G} is interpreted by its
holomorphic extension at one. Local uniform convergence, together with the
classical Gamma product, gives
\[
 \mathcal G(1,z;a)=\frac{\Gamma(a)}{\Gamma(a+z)}.
\]
We evaluate $F_R$ at the fixed argument $z=-a-m$. This is different from
differentiating along the zero path $z=-(a+m)^s$, where the function is
identically zero.

\subsection{A finite primitive coordinate system}

For $x\ge0$, integers $r\ge0$, and integers $k\ge2$, define
\begin{align}
 L_r(x)&=\sum_{n=1}^{\infty}
   \frac{\log^r(n+x)-\log^r n}{n},\qquad L_0(x)=0,
       \label{mz:L}\\
 M_{r,k}(x)&=\sum_{n=1}^{\infty}
          \frac{\log^r(n+x)}{n^k}.
       \label{mz:M}
\end{align}
The convention for the zeroth power is one. The sums and their fixed
derivatives converge locally uniformly: the summands of $L_r$, for $r>0$,
are $O(\log^{r-1}n/n^2)$ on a compact $x$-interval, and those of $M_{r,k}$
are $O(\log^r n/n^k)$. They also define holomorphic functions on
$\mathbb C\setminus(-\infty,-1]$ using principal logarithms.

Put
\begin{equation}
 T_r(x)=L_r(x)+\gamma_r-\gamma_r(1+x).
 \label{mz:T}
\end{equation}

\begin{theorem}[Finite Stieltjes primitive closure]
\label{mz:thm:primitive}
For $r\ge1$ and $x>0$,
\begin{align}
 L_r'(x)&=\frac r{x}T_{r-1}(x),\qquad L_r(0)=0,
      \label{mz:Lrec}\\
 M_{r,k}'(x)
 &=r\left\{\sum_{j=2}^{k}
     \frac{(-1)^{k-j}}{x^{k-j+1}}M_{r-1,j}(x)
       +\frac{(-1)^{k-1}}{x^k}T_{r-1}(x)\right\},
      \label{mz:Mrec}\\
 M_{r,k}(0)&=(-1)^r\zeta^{(r)}(k),
       \qquad M_{0,k}(x)=\zeta(k).
      \label{mz:Minitial}
\end{align}
The combined right sides have removable singularities at zero.
Consequently, any finite rectangle of $L_r$ and $M_{r,k}$ can be
constructed by finitely many ordinary integrations from zero, starting
with ordinary zeta derivatives and the generalized Stieltjes functions.
\end{theorem}

\begin{proof}
The defining limit for generalized Stieltjes constants gives
\begin{align}
 \sum_{n=1}^{\infty}\log^r(n+x)
          \left(\frac1n-\frac1{n+x}\right)
 &=L_r(x)+\gamma_r-\gamma_r(1+x)\notag\\
 &=T_r(x).
 \label{mz:Tsum}
\end{align}
To see the normalization explicitly, subtract the two cutoff sums for
$\gamma_r$ and $\gamma_r(1+x)$. Their integral counterterms differ by
$[\log^{r+1}(N+1+x)-\log^{r+1}(N+1)]/(r+1)$, which tends to zero.
The remaining difference is absolutely convergent. Differentiating
\eqref{mz:L} termwise and using
$1/[n(n+x)]=x^{-1}(1/n-1/(n+x))$ proves \eqref{mz:Lrec}.

For $k\ge2$, the elementary partial-fraction identity
\begin{equation}
 \frac1{n^k(n+x)}
 =\sum_{j=2}^{k}\frac{(-1)^{k-j}}{x^{k-j+1}n^j}
   +\frac{(-1)^{k-1}}{x^k}
          \left(\frac1n-\frac1{n+x}\right)
 \label{mz:partial}
\end{equation}
and termwise differentiation of \eqref{mz:M} prove \eqref{mz:Mrec}.
At zero, differentiation of the absolutely convergent Dirichlet series
for $\zeta(k)$ gives \eqref{mz:Minitial}. The original differentiated
series are holomorphic near zero. Thus the combined rational expressions
on the right of \eqref{mz:Lrec}--\eqref{mz:Mrec} have removable
singularities there. Integrating those combined expressions proves the
last assertion by induction on $r$.
\end{proof}

One must integrate the combined right side of \eqref{mz:Mrec}. Its
individual displayed terms can have nonintegrable poles at zero, although
their sum is regular.

In particular, the first coordinate is the ordinary digamma primitive
\begin{equation}
 D(x):=L_1(x)
 =\sum_{n\ge1}\frac{\log(1+x/n)}n
 =\int_0^x\frac{\psi(1+t)+\gamma}{t}\,dt.
 \label{mz:D}
\end{equation}
The integrand takes the limiting value $\zeta(2)$ at zero. Higher $L_r$
are obtained from
\[
 L_r(x)=r\int_0^x
       \frac{L_{r-1}(t)+\gamma_{r-1}-\gamma_{r-1}(1+t)}t\,dt.
\]
For example
\begin{align}
 L_2'(x)&=\frac2x\bigl(D(x)+\gamma_1-\gamma_1(1+x)\bigr),\\
 M_{2,2}'(x)
 &=\frac2xM_{1,2}(x)
   -\frac2{x^2}\bigl(D(x)+\gamma_1-\gamma_1(1+x)\bigr).
\end{align}

\subsection{The zero-factor theorem}

Let $\left\{\begin{smallmatrix}r\\k\end{smallmatrix}\right\}$ be a
Stirling number of the second kind and set
\begin{equation}
 A_{r,k}=(k-1)!
        \left\{\begin{matrix}r\\k\end{matrix}\right\}.
 \label{mz:A}
\end{equation}
Let $B_j$ be the Bernoulli numbers with $B_1=-1/2$. Define the complete
exponential Bell polynomials by
\[
 \exp\left(\sum_{j\ge1}c_j\frac{u^j}{j!}\right)
       =\sum_{n\ge0}Y_n(c_1,\ldots,c_n)\frac{u^n}{n!}.
\]

\begin{theorem}[All fixed-argument spectral jets at a Gamma zero]
\label{mz:thm:zero}
Let $a>0$, $m\in\mathbb Z_{\ge0}$, and set
\[
 x=a+m,\qquad \ell=\log x,\qquad
 h_0=x\Gamma(a)(-1)^m m!.
\]
For $r\ge1$ define the finite expression
\begin{align}
 b_r={}&(-1)^{r+1}\left\{
      x\bigl(\gamma_r+L_r(x)\bigr)
       +\sum_{k=2}^{r}A_{r,k}x^kM_{r,k}(x)\right.\notag\\
 &\hspace{25mm}\left.
       +\sum_{j=1}^{m}\log^r(x-j)
            \sum_{k=1}^{r}A_{r,k}\left(-\frac xj\right)^k
                        \right\},\label{mz:b}\\
 c_r={}&b_r+\frac{B_r}{r}\ell^r.
 \label{mz:c}
\end{align}
Empty sums vanish. Then, for every $R\ge1$,
\begin{equation}
 \boxed{F_R(-x;a)
    =R h_0\ell\,Y_{R-1}(c_1,\ldots,c_{R-1}).}
 \label{mz:allzero}
\end{equation}
At $x=1$, every $F_R(-1;1)$ is zero. For $x\ne1$, the formula requires
only the $R-1$ logarithmic jets $b_1,\ldots,b_{R-1}$.
\end{theorem}

\begin{proof}
Remove the unique vanishing factor:
\[
 \mathcal G(1+u,-x;a)=(1-e^{-\ell u})\mathcal H(1+u),
\]
where $\mathcal H$ is holomorphic and nonzero near one. At $s=1$, the
Gamma slice has derivative $\Gamma(a)(-1)^m m!$ in $z$ at $-x$.
The removed factor is $(z+x)/x$ on that slice, so $\mathcal H(1)=h_0$.
Equivalently, for a cutoff $N>m$, the zero-removed value is
\begin{equation}
 \mathcal H_N(1)=h_0 (N+a)^x
                      \frac{\Gamma(N-m)}{\Gamma(N+a)},
 \label{mz:Hcut}
\end{equation}
which tends to $h_0$ by the Gamma ratio asymptotic.

The source cutoff is
\[
 \mathcal G_N(s,-x;a)
   =e^{xC_N(s;a)}\prod_{n<N}(1-xx_n^{-s}),\qquad
 C_N(s;a)=\frac{1-(N+a)^{1-s}}{s-1}.
\]
Its derivatives satisfy
$\partial_s^r C_N(1;a)=(-1)^r\log^{r+1}(N+a)/(r+1)$.
For integer $r\ge1$,
\begin{equation}
 \left.\partial_s^r\log(1-xx_n^{-s})\right|_{s=1}
       =(-1)^{r+1}\log^r x_n\,
                         \operatorname{Li}_{1-r}(x/x_n).
 \label{mz:LiDerivative}
\end{equation}
Only integer nonpositive polylogarithm orders occur here, hence these are
rational functions; the formula is valid even for $n<m$, when the
argument exceeds one. The rational identity
\begin{equation}
 \operatorname{Li}_{1-r}(q)
       =\sum_{k=1}^{r}A_{r,k}\left(\frac q{1-q}\right)^k
 \label{mz:LiStirling}
\end{equation}
follows either by applying $q\partial_q$ repeatedly to $q/(1-q)$ or
from the Stirling recurrence. Thus
\begin{align}
 \left.\partial_s^r\log\mathcal H_N(s)\right|_{s=1}
 =(-1)^r\left\{\frac{x}{r+1}\log^{r+1}(N+a)
  -\sum_{\substack{0\le n<N\\n\ne m}}\log^r(n+a)
          \sum_{k=1}^{r}\frac{A_{r,k}x^k}{(n-m)^k}\right\}.
 \label{mz:bcut}
\end{align}
The $n<m$ terms become the finite $j$-sum in \eqref{mz:b}. Set $n=m+j$
in the remaining terms. Since $A_{r,1}=1$,
\[
 \lim_{J\to\infty}\left\{
   \sum_{j=1}^{J}\frac{\log^r(j+x)}j
       -\frac{\log^{r+1}(J+1+x)}{r+1}\right\}
       =\gamma_r+L_r(x).
\]
Indeed, subtract the analogous limit for $\gamma_r$; the change in its
logarithmic counterterm tends to zero. The sums for $k\ge2$ converge to
$M_{r,k}(x)$. Local uniform convergence of the zero-removed products
justifies passage to their logarithmic derivatives, proving
$b_r=\partial_s^r\log\mathcal H(1)$ and \eqref{mz:b}.

For $\ell\ne0$, the Bernoulli generating function yields, near $u=0$,
\begin{equation}
 \log\frac{1-e^{-\ell u}}{\ell u}
       =\sum_{r\ge1}\frac{B_r\ell^r}{r}\frac{u^r}{r!}.
 \label{mz:Bernoulli}
\end{equation}
Consequently
\[
 \mathcal G(1+u,-x;a)
       =h_0\ell u\exp\left(\sum_{r\ge1}c_r\frac{u^r}{r!}\right).
\]
Bell coefficient extraction proves \eqref{mz:allzero}. If $\ell=0$,
the removed factor $1-e^{-\ell u}$ vanishes identically, proving the
stated exceptional case without division by zero.
\end{proof}

The first three nontrivial rows are
\begin{align}
 F_1(-x;a)&=h_0\ell,\label{mz:F1}\\
 F_2(-x;a)&=h_0(2\ell b_1-\ell^2),\label{mz:F2b}\\
 F_3(-x;a)&=h_0\{3\ell(b_1^2+b_2)-3\ell^2b_1+\ell^3\}.
 \label{mz:F3}
\end{align}
The first is the source's already proved zero-flow formula. The second has
the explicit Gamma--Stieltjes--digamma-primitive evaluation
\begin{equation}
 \boxed{F_2(-x;a)=h_0\left[
   2x\ell\left\{\gamma_1+D(x)
           -\sum_{j=1}^{m}\frac{\log(x-j)}j\right\}-\ell^2\right].}
 \label{mz:F2}
\end{equation}
For the next row,
\begin{align}
 b_2=-\biggl[&x(\gamma_2+L_2(x))+x^2M_{2,2}(x)\notag\\
 &+\sum_{j=1}^{m}\log^2(x-j)
                    \left(-\frac xj+\frac{x^2}{j^2}\right)\biggr].
 \label{mz:b2}
\end{align}

For fixed $R\ge2$, Theorem~\ref{mz:thm:primitive} supplies a finite construction
of all entries needed in \eqref{mz:allzero}. It uses ordinary Stieltjes
constants only through $\gamma_{R-1}$, Stieltjes \emph{functions}
$\gamma_j(1+t)$ only through $j=R-2$, and finitely many zeta derivatives
$\zeta^{(r)}(k)$ with $0\le r\le R-1$ and $2\le k\le R-1$.
For $R\ge2$, the highest ordinary Stieltjes constant occurs linearly,
with coefficient
\begin{equation}
 [\gamma_{R-1}]F_R(-x;a)
       =(-1)^R R x h_0\log x
 \label{mz:highest}
\end{equation}
in this explicit coordinate formula. This is a statement about the
displayed construction. It does not assert arithmetic independence of
Stieltjes constants from zeta derivatives or their integrals.

\subsection{Polylogarithmic Mellin realization}

The coordinate system also has a direct polylogarithm interpretation. Set
\begin{equation}
 J_k(v,x)=\sum_{n=1}^{\infty}\frac1{n^k(n+x)^v}.
 \label{mz:J}
\end{equation}
For $x\ge0$ and $\Re v>0$, termwise integration of the absolutely
convergent polylogarithm series gives
\begin{equation}
 J_k(v,x)=\frac1{\Gamma(v)}\int_0^{\infty}
     t^{v-1}e^{-xt}\operatorname{Li}_k(e^{-t})\,dt,
       \qquad k\ge1.
 \label{mz:Mellin}
\end{equation}
For $k\ge2$ its Dirichlet series is holomorphic near $v=0$, and
\[
 M_{r,k}(x)=(-1)^r\partial_v^rJ_k(0,x).
\]
For $k=1$, the difference
\[
 J_1(v,x)-\zeta(1+v)
 =\sum_{n\ge1}\frac{(n+x)^{-v}-n^{-v}}n
\]
is holomorphic near zero because its summands are $O(n^{-2-\Re v})$
there, uniformly on compact sets. It follows that
\begin{equation}
 J_1(v,x)=\frac1v+
       \sum_{r\ge0}\frac{(-1)^r(\gamma_r+L_r(x))}{r!}v^r.
 \label{mz:JLaurent}
\end{equation}
Thus the exact Gamma-zero derivatives are finite combinations of resonant
Laurent coefficients and regular spectral derivatives of the elementary
polylogarithmic Mellin transform \eqref{mz:Mellin}. This is a specific
shifted Dirichlet series, not an identification of its Laurent coefficients
with the ordinary Hurwitz Stieltjes constants.

\subsection{Generalized harmonic numbers give every coefficient}

Let $e_j(h)$ be defined by
\begin{equation}
 \prod_{n=1}^{h}(1+u/n)=\sum_{j=0}^{h}e_j(h)u^j,
 \qquad e_0(h)=1.
 \label{mz:e}
\end{equation}
These are polynomials in the generalized harmonic numbers
$H_h^{(j)}=\sum_{n=1}^h n^{-j}$:
\[
 e_1(h)=H_h,\qquad
 e_2(h)=\tfrac12\bigl(H_h^2-H_h^{(2)}\bigr),\qquad
 \sum_j e_j(h)u^j
  =\exp\left(\sum_{j\ge1}\frac{(-1)^{j-1}H_h^{(j)}}j u^j\right).
\]
For $r\ge1$, $h\ge1$, and $k\ge1$, put
\begin{equation}
 Q_{r,k,h}(q)
   =\sum_{j=1}^{\min(r,h)}\frac{r!}{(r-j)!}
           e_{j-1}(h-1)\zeta^{(r-j)}(k+h,q).
 \label{mz:Q}
\end{equation}

\begin{theorem}[Harmonic--Hurwitz coefficient formulas]
\label{mz:thm:coefficients}
Let $r\ge1$, let $N\ge0$ be an integer, and let
$x\in\mathbb C\setminus(-\infty,-1]$ satisfy $|x|<N+1$. Then
\begin{align}
 L_r(x)={}&\sum_{n=1}^{N}
        \frac{\log^r(n+x)-\log^r n}{n}
   +(-1)^r\sum_{h\ge1}\frac{(-x)^h}{h}Q_{r,1,h}(N+1),
       \label{mz:Lseries}\\
 M_{r,k}(x)={}&\sum_{n=1}^{N}\frac{\log^r(n+x)}{n^k}
       +(-1)^r\zeta^{(r)}(k,N+1)\notag\\
 &+(-1)^r\sum_{h\ge1}\frac{(-x)^h}{h}Q_{r,k,h}(N+1),
       \qquad k\ge2.
       \label{mz:Mseries}
\end{align}
The series converge normally on compact subsets of $|x|<N+1$.
Every positive $x$ is covered after a finite cutoff. The only harmonic
orders needed for a spectral derivative of order $r$ are at most $r-1$.
\end{theorem}

\begin{proof}
Split \eqref{mz:J} at $N$ and use the binomial expansion in the tail:
\[
 \sum_{n>N}\frac1{n^k(n+x)^v}
   =\sum_{h\ge0}\frac{(-1)^h(v)_h x^h}{h!}
                                      \zeta(k+v+h,N+1).
\]
For $|x|<N+1$, the double series and its fixed $v$ derivatives converge
normally on the appropriate local domain; the $h=0,k=1$ pole is treated
separately. Since
\[
 \frac{(v)_h}{h!}=\frac vh
       \prod_{j=1}^{h-1}(1+v/j)
       =\frac1h\sum_{j=1}^{h}e_{j-1}(h-1)v^j,
\]
Leibniz differentiation at zero gives the coefficient
$Q_{r,k,h}(N+1)/h$. Formula \eqref{mz:JLaurent} or the derivative
description of $M_{r,k}$ then proves \eqref{mz:Lseries}--\eqref{mz:Mseries}.
The finite harmonic expression follows from \eqref{mz:e}.
\end{proof}

For example, for $|x|<1$,
\begin{align}
 D(x)&=\sum_{h\ge1}\frac{(-1)^{h+1}x^h}{h}\zeta(h+1),\\
 L_2(x)&=2\sum_{h\ge1}\frac{(-x)^h}{h}
       \bigl\{\zeta'(h+1)+H_{h-1}\zeta(h+1)\bigr\},\\
 M_{2,2}(x)&=\zeta''(2)+2\sum_{h\ge1}\frac{(-x)^h}{h}
       \bigl\{\zeta'(h+2)+H_{h-1}\zeta(h+2)\bigr\}.
\end{align}

\paragraph{Precise remaining questions.}
The finite coordinate construction is complete, but whether some of its
nontrivial fixed-argument integrals reduce to a smaller field of special
values remains open. The first test is $D(x)$ at a positive integer or a
nonintegral rational point. No failure of a symbolic or numerical search
would prove non-reducibility. A second target is a global description of
the individual higher spectral jets as entire functions of $z$, with these
Gamma-zero evaluations as interpolation data. A third target is the
compatible construction for several simultaneously moving zero factors
in multicolour spectral products.

% END sections/05-gamma-roots.tex
% BEGIN sections/06-shifted-polylogarithms.tex
% Insertion-ready research notes; uses standard amsmath/amsthm notation.
% Derivation and numerical audit: 2026-10-10.
% Do not advertise a new general parity principle: see the scope paragraph.

\section{A two-shift Lerch reflection for inverse-color doubles}
\label{gap:sec:main}\label{sec:gap}

The manuscript already proves an all-weight reduction of
$\operatorname{Li}_{a,b}(z,z^{-1})$ to one-direction double
polylogarithms. Its explicit mixed parity formula treats $b=1$; the
general parity component also follows by combining that reduction with
the all-index one-direction parity theorem. Thus an extension from
$b=1$ alone must not be advertised as a new reducibility principle.
The established multiple-polylogarithm parity theory includes Panzer's
general theorem and explicit depth-two formulas \cite{panzer}. More recently, Xu's
2026 preprint \emph{Symmetry Results for Cyclotomic Multiple Hurwitz
Zeta Values via Contour Integrals}, arXiv:2602.10391 \cite{xu}, proves broad
shifted cyclotomic symmetry identities. Its index order is increasing,
whereas the manuscript uses decreasing indices.

The purpose of the following calculation is more specific: it gives a
compact two-shift identity that is an ordinary holomorphic identity in
the full disk of a continuous color variable. Its arbitrary-index
version follows from a single first-order seed by parameter
differentiation. It also gives complete reductions, rather than only
paired reductions, on the complementary-shift diagonal. These are
useful explicit additions to the manuscript. No claim of worldwide
priority is made for the underlying symmetry.

\subsection{Definitions and analytic domain}

For integers $a,b\geq1$, define
\begin{equation}\label{gap:eq:def}
 \begin{aligned}
 \mathcal M_{a,b}(z;u,v)
 &=\sum_{n>m\geq0}
   \frac{z^{n-m}}{(n+u)^a(m+v)^b}
 =\sum_{h\geq1}z^h c_{a,b}(h;u,v),\\
 c_{a,b}(h;u,v)&=\sum_{m\geq0}
       \frac1{(m+h+u)^a(m+v)^b}.
 \end{aligned}
\end{equation}
We initially work on
\begin{equation}\label{gap:eq:domain}
 |z|<1,\qquad 0<\operatorname{Re}u<1,
 \qquad 0<\operatorname{Re}v<1.
\end{equation}
All powers are integer powers, so there is no denominator branch
choice. The Lerch function is
\[
 \Phi(z,s,q)=\sum_{m\geq0}\frac{z^m}{(m+q)^s}
 \quad (|z|<1,\ \operatorname{Re}q>0).
\]
Put
\begin{equation}\label{gap:eq:C}
 C_1(q)=\pi\cot\pi q,
 \qquad
 C_r(q)=\zeta(r,q)+(-1)^r\zeta(r,1-q)\quad(r\geq2).
\end{equation}
Equivalently,
\begin{equation}\label{gap:eq:Cderivative}
 C_r(q)=\frac{(-1)^{r-1}}{(r-1)!}
          \frac{d^{r-1}}{dq^{r-1}}(\pi\cot\pi q).
\end{equation}
Thus every $C_r$ is elementary in trigonometric functions, and can
equally be expressed through polygamma functions.

The series in \eqref{gap:eq:def} converges absolutely and normally on
compact subsets of \eqref{gap:eq:domain}. Indeed, uniformly on compact
shift sets,
\[
 |c_{a,b}(h;u,v)|\leq
 C\sum_{m\geq0}\frac1{(m+h+1)^a(m+1)^b}
 =O\left(\frac{1+\log(h+1)}{(h+1)^a}\right).
\]
The same estimates hold after any fixed number of shift derivatives.
In particular all parameter differentiations below are legitimate
holomorphic differentiations, not formal manipulations of a divergent
series.

\subsection{The first-order seed}

\begin{theorem}[Two-shift seed]\label{gap:thm:seed}
On \eqref{gap:eq:domain},
\begin{align}
 &\mathcal M_{1,1}(z;u,v)
   +\mathcal M_{1,1}(z;1-v,1-u)\notag\\
 &\quad=z\Phi(z,1,u)\Phi(z,1,1-v)
   +\bigl(C_1(v)-C_1(u)\bigr)
       z\Phi(z,1,1+u-v).
 \label{gap:eq:seed}
\end{align}
All apparent confluences at $u=v$ are removable; in fact no divided
difference is needed in the displayed formula.
\end{theorem}

\begin{proof}
Fix a positive integer $h$. The absolutely convergent bilateral sum
\[
 B_{1,1}(h;u,v)=\sum_{m\in\mathbb Z}
       \frac1{(m+h+u)(m+v)}
\]
has three disjoint ranges. The range $m\geq0$ is $c_{1,1}(h;u,v)$.
In the range $m\leq-h-1$, set $m=-h-1-\ell$; its contribution is
$c_{1,1}(h;1-v,1-u)$. In the middle range $-h\leq m\leq-1$, write
$m=-\ell-1$ and $j=h-1-\ell$. Its contribution is
\[
 -\sum_{j+\ell=h-1}\frac1{(j+u)(\ell+1-v)}.
\]
On the other hand, with $d=h+u-v$,
\[
 \frac1{(m+h+u)(m+v)}
 =\frac1d\left(\frac1{m+v}-\frac1{m+h+u}\right).
\]
Symmetric summation gives
$B_{1,1}(h;u,v)=(C_1(v)-C_1(u))/(h+u-v)$.
The shift by the integer $h$ does not change a symmetric cotangent
sum: the finitely shifted endpoint terms tend to zero. Multiply by
$z^h$ and sum over $h\geq1$. The middle range is the coefficient
convolution of the two Lerch series, and the bilateral term sums to
the second term in \eqref{gap:eq:seed}. All resulting $h$-series
converge absolutely for $|z|<1$.
\end{proof}

\subsection{Every pair of indices, in two equivalent forms}

\begin{theorem}[Arbitrary-index reflection]\label{gap:thm:all}
Let $a,b\geq1$ and $w=a+b$. Set $q=1+u-v$. Then
\begin{align}
 &\mathcal M_{a,b}(z;u,v)
      +(-1)^w\mathcal M_{b,a}(z;1-v,1-u)\notag\\
 &=z\Bigg[
 (-1)^b\sum_{r=1}^{a}\binom{w-r-1}{b-1}
       C_r(u)\Phi(z,w-r,q)\notag\\
 &\hspace{17mm}
 +\sum_{r=1}^{b}(-1)^{b-r}\binom{w-r-1}{a-1}
       C_r(v)\Phi(z,w-r,q)\notag\\
 &\hspace{17mm}
 -(-1)^b\Phi(z,a,u)\Phi(z,b,1-v)\Bigg].
 \label{gap:eq:all}
\end{align}
Equivalently, the right side is
\begin{equation}\label{gap:eq:all-derivative}
 \frac{(-1)^w}{(a-1)!(b-1)!}
 \partial_u^{a-1}\partial_v^{b-1}
 \left[
 z\Phi(z,1,u)\Phi(z,1,1-v)
 +(C_1(v)-C_1(u))z\Phi(z,1,1+u-v)
 \right].
\end{equation}
\end{theorem}

\begin{proof}
For the derivative form, apply the stated mixed derivative to
\eqref{gap:eq:seed}. On the first summand on its left, the derivative
contributes the sign $(-1)^{a+b-2}=(-1)^w$ and the factorial factor.
On the reflected summand both chain-rule signs cancel the signs
from differentiating a reciprocal, and the two indices are exchanged.
This gives \eqref{gap:eq:all-derivative} with the precise sign shown.

For an independent coefficient proof of \eqref{gap:eq:all}, put
$d=h+u-v$. The rational partial-fraction identity is
\begin{align}
 \frac1{(m+h+u)^a(m+v)^b}
 ={}&(-1)^b\sum_{r=1}^{a}
  \binom{w-r-1}{b-1}\frac1{d^{w-r}(m+h+u)^r}\notag\\
 &+\sum_{r=1}^{b}(-1)^{b-r}
  \binom{w-r-1}{a-1}\frac1{d^{w-r}(m+v)^r}.
 \label{gap:eq:pf}
\end{align}
It follows either by expanding at its two poles or by subtracting
the two sides, observing that the difference has no finite poles and
vanishes at infinity. Since $\operatorname{Re}d>0$, no collision occurs
in the domain under consideration.

Sum over $m\in\mathbb Z$. For $r\geq2$, the sum of each reciprocal
power is $C_r(u)$ or $C_r(v)$. The $r=1$ terms are evaluated by a
common symmetric cutoff. Their principal parts at infinity cancel,
because
\[
 \binom{w-2}{a-1}=\binom{w-2}{b-1}.
\]
The original bilateral sum is absolutely convergent, so this use of
a common cutoff changes no value. Split it into the same three ranges
as in the first-order proof. The negative range is now
$(-1)^w c_{b,a}(h;1-v,1-u)$, and the middle range is
\[
 (-1)^b\sum_{j+\ell=h-1}
      (j+u)^{-a}(\ell+1-v)^{-b}.
\]
Multiplication by $z^h$ and summation gives
\eqref{gap:eq:all}.
\end{proof}

\subsection{Shift derivatives and normalized primitives}

The whole family is closed under arbitrary shift derivatives:
\begin{equation}\label{gap:eq:shift-jets}
 \partial_u^r\partial_v^s\mathcal M_{a,b}(z;u,v)
 =(-1)^{r+s}(a)_r(b)_s\mathcal M_{a+r,b+s}(z;u,v).
\end{equation}
Accordingly, for $a\geq2$ a primitive in $u$ is
$-\mathcal M_{a-1,b}/(a-1)$, and for $b\geq2$ a primitive in $v$ is
$-\mathcal M_{a,b-1}/(b-1)$. At the index-one boundary one should
normalize the primitive before summing. For a path between $u_0$ and
$u_1$ inside the right half-plane,
\begin{equation}\label{gap:eq:log-primitive}
 \int_{u_0}^{u_1}\mathcal M_{1,b}(z;u,v)\,du
 =\sum_{n>m\geq0}\frac{z^{n-m}}{(m+v)^b}
       \log\frac{n+u_1}{n+u_0}.
\end{equation}
Here the logarithm is the difference of the common analytic logarithms
along the path. For $|z|<1$ the normalized series is absolutely
convergent, including $b=1$, because the logarithmic difference is
$O(n^{-1})$ on compact endpoint sets. The unnormalized series with
$\log(n+u)$ alone need not converge when $b=1$.

There is also a precise spectral-derivative interpretation. Define the
termwise difference
\[
 \Delta\mathcal M_{s,b}
 =\sum_{n>m\geq0}\frac{z^{n-m}}{(m+v)^b}
       \bigl((n+u_1)^{-s}-(n+u_0)^{-s}\bigr).
\]
It is holomorphic near $s=0$, even when the two separate sums there
are divergent, and the right side of \eqref{gap:eq:log-primitive} is
$-\partial_s\Delta\mathcal M_{s,b}|_{s=0}$. This statement uses
subtraction before continuation. It does not assert that the
integer-index finite sum \eqref{gap:eq:all} can itself be
differentiated in a noninteger exponent.

\paragraph{Boundary values.}
The identity extends to each $|z|=1$, $z\ne1$, as an ordinary
convergent identity. To make this statement independent of an
unspecified Abel prescription, note that on compact shift sets
$c_{a,b}(h;u,v)\to0$ and
\[
 c_{a,b}(h+1;u,v)-c_{a,b}(h;u,v)
 =O\left(\frac{1+\log(h+1)}{(h+1)^{a+1}}\right).
\]
The sum of these differences converges absolutely. Summation by
parts therefore proves convergence uniformly on closed arcs that
avoid $1$. The same argument applies to every Lerch term on the
right, whose parameter has positive real part, and to fixed shift
derivatives. Taking a radial limit of the disk identity proves the
boundary identity in the ordinary gap-ordered series convention.

It also agrees with the original outer-index order. Let $O_N$ be the
sum in \eqref{gap:eq:def} restricted to $n\leq N$, and let
$D_N=\sum_{h=1}^Nz^hc_{a,b}(h;u,v)$. Partial summation of the finite
$h$-sum gives, uniformly for shifts in a compact subset of the strip,
\[
 |D_N-O_N|\leq\frac{C}{|1-z|}
 \left(N^{-a}(1+\log(N+1))+\sum_{m>N}m^{-a-b}\right)
 \longrightarrow0.
\]
For $m\leq N$, the omitted $h$-range starts past $N-m$, so its
denominator $(m+h+u)^a$ is of size at least $N^a$; for $m>N$, the
full finite $h$-sum is bounded by a constant times $m^{-a}$.
These are exactly the two terms in the displayed bound. Thus no
conditional rearrangement is hidden in the boundary statement.
If $a,b\geq2$,
all terms are also absolutely convergent at $z=1$, so that endpoint
may then be included.

\subsection{Complete reductions at complementary shifts}

The reflection becomes a full evaluation when the two exponents
agree and $v=1-u$.

\begin{corollary}[Complementary diagonal]\label{gap:cor:diagonal}
For $a\geq1$, $|z|<1$ and $0<\operatorname{Re}u<1$,
\begin{align}
 \mathcal M_{a,a}(z;u,1-u)
 =(-1)^a z\Bigg[
 &\sum_{r=1}^{a}\binom{2a-r-1}{a-1}
        C_r(u)\Phi(z,2a-r,2u)\notag\\
 &-\frac12\Phi(z,a,u)^2\Bigg].
 \label{gap:eq:diagonal}
\end{align}
In particular,
\begin{equation}\label{gap:eq:diagonal11}
 \mathcal M_{1,1}(z;u,1-u)
 =\frac z2\Phi(z,1,u)^2
   -\pi\cot(\pi u)\,z\Phi(z,1,2u),
\end{equation}
and
\begin{align}
 \mathcal M_{2,2}(z;u,1-u)
 =z\bigl[&2\pi\cot(\pi u)\Phi(z,3,2u)
          +\pi^2\csc^2(\pi u)\Phi(z,2,2u)\notag\\
          &-\tfrac12\Phi(z,2,u)^2\bigr].
 \label{gap:eq:diagonal22}
\end{align}
\end{corollary}
\begin{proof}
The reflected double on the left of \eqref{gap:eq:all} equals the
original one. On the right, use $C_r(1-u)=(-1)^rC_r(u)$ and divide
by two.
\end{proof}

The $a=1$ identity also exposes the generalized-harmonic content.
With
\[
 H_h(u)=\psi(h+u)-\psi(u)=\sum_{j=0}^{h-1}\frac1{j+u},
\]
partial fractions give
\[
 c_{1,1}(h;u,1-u)
 =\frac{H_h(u)-\pi\cot\pi u}{h+2u-1}.
\]
Consequently \eqref{gap:eq:diagonal11} is equivalent to
\begin{equation}\label{gap:eq:harmonic-convolution}
 \sum_{h\geq1}\frac{z^h H_h(u)}{h+2u-1}
       =\frac z2\Phi(z,1,u)^2.
\end{equation}
This last elementary convolution identity can be proved directly by
partial fractions in the coefficient of $z^{h-1}$ in the square.
That observation clarifies the strength of the claim: the result is
an exact complete reduction, but the first-order instance already has
an elementary coefficient explanation.

\subsection{The half-shift diagonal in classical polylogarithms}

Put
\[
 \chi_s(x)=\frac{\operatorname{Li}_s(x)-\operatorname{Li}_s(-x)}2.
\]
For $z=x^2$ with $|x|<1$,
\[
 \Phi(z,s,\tfrac12)=\frac{2^s}{x}\chi_s(x),\qquad
 C_{2r}(\tfrac12)=2(2^{2r}-1)\zeta(2r),\qquad
 C_{2r+1}(\tfrac12)=0.
\]
The dependence on the chosen square root cancels in all products.
Equation \eqref{gap:eq:diagonal} gives the all-order identity
\begin{align}
 \mathcal M_{a,a}(z;\tfrac12,\tfrac12)
 =(-1)^a\Bigg[
 &2\sum_{r=1}^{\lfloor a/2\rfloor}
   (2^{2r}-1)\binom{2a-2r-1}{a-1}\zeta(2r)
                         \operatorname{Li}_{2a-2r}(z)\notag\\
 &-2^{2a-1}\chi_a(\sqrt z)^2\Bigg].
 \label{gap:eq:half-all}
\end{align}
The first four cases are
\begin{align}
 \mathcal M_{1,1}(z;\tfrac12,\tfrac12)
   &=2\operatorname{artanh}^2\sqrt z,\label{gap:eq:half11}\\
 \mathcal M_{2,2}(z;\tfrac12,\tfrac12)
   &=\pi^2\operatorname{Li}_2(z)-8\chi_2(\sqrt z)^2,
      \label{gap:eq:half22}\\
 \mathcal M_{3,3}(z;\tfrac12,\tfrac12)
   &=32\chi_3(\sqrt z)^2-3\pi^2\operatorname{Li}_4(z),
      \label{gap:eq:half33}\\
 \mathcal M_{4,4}(z;\tfrac12,\tfrac12)
   &=10\pi^2\operatorname{Li}_6(z)
       +\frac{\pi^4}{3}\operatorname{Li}_4(z)
       -128\chi_4(\sqrt z)^2.
      \label{gap:eq:half44}
\end{align}
Here each left side is a two-index series at a variable
argument; the displayed equality is a reduction of that function.
At individual special points there may be additional elementary
stuffle explanations, which should be recorded rather than presented
as new arithmetic phenomena.

\subsection{A classical antiderivative and an exact alternating moment}

At the half shift,
\[
 H_h(\tfrac12)=2H_{2h}-H_h.
\]
Hence \eqref{gap:eq:half11} says
\begin{equation}\label{gap:eq:harmonic-half}
 \sum_{h\geq1}\frac{(2H_{2h}-H_h)z^h}{h}
     =2\operatorname{artanh}^2\sqrt z.
\end{equation}
Integrating once against $dz/z$ is a useful explicit antiderivative.
For $0<z<1$ put \mbox{$x=(1-\sqrt z)/(1+\sqrt z)$}. Then
\begin{align}
 \sum_{h\geq1}\frac{(2H_{2h}-H_h)z^h}{h^2}
 ={}&\frac72\zeta(3)
   -\log^2x\log\frac{1+x}{1-x}\notag\\
   &+4\log x\,\chi_2(x)-4\chi_3(x).
 \label{gap:eq:primitive}
\end{align}
To prove it, write the left side as
$4\int_0^{\sqrt z}\operatorname{artanh}^2t\,dt/t$ and substitute
$x=(1-t)/(1+t)$. An antiderivative of
$\log^2x((1-x)^{-1}+(1+x)^{-1})$ is
\[
 \log^2x\log\frac{1+x}{1-x}
 -4\log x\,\chi_2(x)+4\chi_3(x).
\]
At $x=1$ its value is $4\chi_3(1)=7\zeta(3)/2$, fixing the constant.
The real interval formulation avoids any hidden logarithm branches;
continuation elsewhere is along the corresponding common path.

At the alternating endpoint one obtains
\begin{equation}\label{gap:eq:alternating-primitive}
 \sum_{h\geq1}\frac{(-1)^h(2H_{2h}-H_h)}{h^2}
       =\frac72\zeta(3)-2\pi G,
 \qquad G=\beta(2).
\end{equation}
An entirely real proof follows from
$-4\int_0^1\arctan^2t\,dt/t$; alternatively continue
\eqref{gap:eq:primitive} to $z=-1$, where $x=-i$ and the imaginary
$\pi^3$ terms cancel. The series is absolutely convergent.
This is a checkable example in the existing Euler-sum vocabulary,
not a claim that this classical type of special value is new.

\subsection{Rational shifts and finite root-of-unity filters}

For integers $1\leq p\leq q$, put $\omega=e^{2\pi i/q}$ and choose
$x$ with $x^q=z$. The residue-class filter gives
\begin{equation}\label{gap:eq:rational-filter}
 \Phi(z,s,p/q)=q^{s-1}x^{-p}
    \sum_{j=0}^{q-1}\omega^{-pj}
                    \operatorname{Li}_s(\omega^j x).
\end{equation}
For integer $s\geq1$ this follows directly by expanding the
polylogarithms and projecting onto exponents congruent to $p$ modulo
$q$. For larger positive rational shifts, first use
\[
 \Phi(z,s,q_0+N)=z^{-N}
 \left(\Phi(z,s,q_0)-\sum_{j=0}^{N-1}\frac{z^j}{(j+q_0)^s}\right).
\]
All apparent singularities at $z=0$ are removable. Thus rational
instances of \eqref{gap:eq:all} and \eqref{gap:eq:diagonal} reduce to
finite expressions involving ordinary polylogarithms at algebraically
related arguments and elementary cotangent derivatives. No extra
depth is introduced by evaluating the single Lerch functions.

\subsection{An independent integral check and reproducibility}

For $a,b\geq1$, $\operatorname{Re}u,\operatorname{Re}v>0$ and
$|z|<1$,
\begin{equation}\label{gap:eq:integral}
 \mathcal M_{a,b}(z;u,v)
 =\frac z{(a-1)!}\int_0^1
       \frac{t^u(-\log t)^{a-1}\Phi(t,b,v)}{1-zt}\,dt.
\end{equation}
Insert the Mellin integral for $(n+u)^{-a}$ and sum the geometric
series in $h=n-m$. This proves the formula by absolute convergence.
At a fixed $|z|=1$, $z\ne1$, its denominator stays away from zero on
$[0,1]$; the possible logarithmic singularity of $\Phi(t,1,v)$ at
$t=1$ is integrable. The boundary formula follows by dominated
convergence and the preceding summation-by-parts argument.

The supplied script \texttt{check\_shifted\_gap.py} evaluates this
integral after $t=e^{-y}$, using finite root-of-unity filters for
rational inner shifts. It compares that independent quadrature with
\eqref{gap:eq:all} for unequal shifts and complex colors, and checks
the first three half-diagonal formulas. At 60 decimal working digits,
the largest observed absolute residual among the four general
reflection checks is $8.33\times10^{-60}$. These checks are numerical
diagnostics, not rigorous error enclosures; the proofs are the
bilateral splitting and partial-fraction arguments above.

\subsection{Precise continuation questions}

\begin{enumerate}
\item Compare \eqref{gap:eq:all} term by term with the depth-two
specialization of Xu's shifted cyclotomic symmetry formulas, recording
the increasing/decreasing index conversion and color normalization.
The existence of a broad earlier theorem should remain explicit.
\item Determine which non-diagonal parameter specializations make the
reflected double equal to, or otherwise independently reducible from,
the original. Complementary shifts with equal exponents give the
complete reduction proved here.
\item Develop closed primitives of the complementary diagonal for
$a\geq2$ in a specified classical or multiple-polylogarithm basis.
The half-shift first-order primitive \eqref{gap:eq:primitive} supplies
a nontrivial normalization check.
\item Study noninteger spectral indices without differentiating the
integer-index partial-fraction formula as though its summation range
were analytic. Such an extension needs a separate contour or Mellin
argument; the integer formula alone does not justify it.
\item Continue in the shifts beyond the safe strip while tracking the
finitely many terms crossed when a shifted lattice point passes zero.
This should be a meromorphic identity with explicit finite
corrections, not an unqualified replacement of ordinary sums by
principal values.
\end{enumerate}

% END sections/06-shifted-polylogarithms.tex
% BEGIN sections/07-spectral-products.tex
\section{Local identities for every derivative of a spectral product}
\label{sec:spectral}

There is a useful second application of the distinction between a
multivariate germ and its one-dimensional restrictions. For polynomially
weighted products of shifted linear factors, every fixed spectral
derivative is a polynomial in the \emph{normalized factor exponents}.
Its highest-degree component is local and explicitly computable. This
produces finite identities between higher zeta derivatives even when the
individual values do not have short closed forms.

The first-derivative multiplicative anomaly is classical in substance;
Dowker treats linear factors and polynomial multiplicities
\cite{dowker}, and G\"urel gives general regularized-product formulas
\cite{gurel}, extending the Lerch--Mizuno setting \cite{mizuno}.
We retain that attribution. The all-order germ, the normalized-weight
polynomial formulation, and its explicit finite-difference consequences
are the proved extension developed here; the literature review does not
establish a priority claim for them.

\subsection{An exact multivariate meromorphic germ}

Let $P$ be a nonzero polynomial of degree $D$, let
$a_1,\ldots,a_m>0$, and initially define
\begin{equation}\label{sp:def}
 \cZ_P(\boldsymbol u;\boldsymbol a)
  =\sum_{n\ge0}P(n)\prod_{j=1}^m(n+a_j)^{-u_j},
 \qquad U=\sum_{j=1}^m u_j.
\end{equation}
Absolute convergence requires $\Rea U>D+1$; no separate positivity of
the $u_j$ is needed. Choose $b>0$ such that
$\max_j|a_j-b|<b$, for example $b=\max_j a_j$. Write
\begin{align}
 d_j&=a_j-b,& p_k&=\frac{P^{(k)}(-b)}{k!},&
 P(n)&=\sum_{k=0}^D p_k(n+b)^k,\label{sp:base}\\
 L_j(z)&=\log(1+d_jz),&
 L_{\boldsymbol u}(z)&=\sum_j u_jL_j(z),&
 C_\ell(\boldsymbol u)&=[z^\ell]e^{-L_{\boldsymbol u}(z)}.
 \label{sp:C}
\end{align}
The logarithms are their analytic germs at $z=0$. Explicitly,
\begin{equation}\label{sp:Cpoch}
 C_\ell(\boldsymbol u)=
 \sum_{q_1+\cdots+q_m=\ell}
       \prod_{j=1}^m\frac{(-1)^{q_j}d_j^{q_j}(u_j)_{q_j}}{q_j!}.
\end{equation}
For $\ell\ge1$ this is a polynomial of degree at most $\ell$ and has
zero constant term.

\begin{theorem}[Finite polar numerator, normally convergent regular part]
\label{sp:thm:germ}
The continuation of \eqref{sp:def} is
\begin{equation}\label{sp:continuation}
 \cZ_P(\boldsymbol u;\boldsymbol a)
   =\sum_{k=0}^D p_k\sum_{\ell\ge0}
          C_\ell(\boldsymbol u)\zeta(U+\ell-k,b).
\end{equation}
Near $\boldsymbol u=0$ it has the exact form
\begin{equation}\label{sp:germ}
 \boxed{\cZ_P(\boldsymbol u;\boldsymbol a)
             =\frac{\cR_P(\boldsymbol u)}U+\cH_P(\boldsymbol u),
 \qquad \cR_P=\sum_{k=0}^D p_k C_{k+1}.}
\end{equation}
Here $\cR_P(0)=0$, $\deg\cR_P\le D+1$, and $\cH_P$ is holomorphic.
An explicit regular part is
\begin{align}\label{sp:H}
 \cH_P(\boldsymbol u)=\sum_{k=0}^D p_k\Bigg\{&\zeta(U-k,b)
   +\sum_{\substack{\ell\ge1\\\ell\ne k+1}}
            C_\ell(\boldsymbol u)\zeta(U+\ell-k,b)\notag\\
   &+C_{k+1}(\boldsymbol u)g_b(U)\Bigg\}.
\end{align}
Every infinite sum displayed here converges normally, with every fixed
spectral derivative, away from its explicitly displayed poles.
\end{theorem}

\begin{proof}
Expand $(n+a_j)^{-u_j}$ as $(n+b)^{-u_j}$ times the binomial series
in $d_j/(n+b)$, then multiply the expansions. Choose a radius $\rho$
with $1/b<\rho<1/\max|d_j|$; when all $d_j=0$, any $\rho>1/b$
will do. On compact $\boldsymbol u$-sets the exponential in
\eqref{sp:C} is bounded on $|z|=\rho$, so
$|C_\ell(\boldsymbol u)|\le M\rho^{-\ell}$.
For sufficiently large $\ell$, the Dirichlet series for
$\zeta(U+\ell-k,b)$ and its fixed derivatives are bounded by a
constant times $b^{-\ell}$, allowing harmless polynomial factors in
$\ell$. Since $b\rho>1$, the resulting series converges geometrically.
The same Cauchy argument on a slightly larger compact set handles
$\boldsymbol u$ derivatives. This proves every interchange in the
initial half-space and then proves \eqref{sp:continuation} by continuation.
For $|U|<1/2$, the only Hurwitz poles occur when $\ell=k+1$.
Subtracting exactly their $1/U$ terms gives \eqref{sp:germ} and
\eqref{sp:H}.
\end{proof}

The germ need not be jointly holomorphic at zero: $\cR_P$ need not
be divisible by $U$. Its restriction to every ray
$\boldsymbol u=s\boldsymbol w$ with $W:=\sum_jw_j\ne0$ is nevertheless
holomorphic at $s=0$, since $\cR_P(0)=0$. The auxiliary $b$ can change
$\cR_P$ off $U=0$ by a multiple of $U$, absorbed into $\cH_P$.
Its restriction to the polar hyperplane, and the tangent coefficient
identities below, are intrinsic.

\subsection{Polynomiality in normalized exponents and exact finite differences}

For $\boldsymbol\alpha$ satisfying $\sum_j\alpha_j=1$, set
\begin{equation}\label{sp:F}
 \mathfrak F_r(\boldsymbol\alpha)
   =\left.\frac{d^r}{ds^r}
          \cZ_P(s\boldsymbol\alpha;\boldsymbol a)\right|_{s=0},
 \qquad r\ge0.
\end{equation}
Complex normalized exponents are allowed by continuation. Write
$\cH_{P,[q]}$ and $\cR_{P,[q]}$ for the homogeneous Taylor parts.

\begin{theorem}[Every normalized spectral jet is a polynomial]
\label{sp:thm:polynomial}
One has
\begin{align}
 \mathfrak F_r(\boldsymbol\alpha)
  &=r!\{\cH_{P,[r]}(\boldsymbol\alpha)
                +\cR_{P,[r+1]}(\boldsymbol\alpha)\},\label{sp:homogeneous}\\
 \cR_{P,[r+1]}(\boldsymbol\alpha)
  &=\frac{(-1)^{r+1}}{(r+1)!}
      \sum_{k=0}^D p_k[z^{k+1}]L_{\boldsymbol\alpha}(z)^{r+1}.
 \label{sp:Rhomogeneous}
\end{align}
Thus $\mathfrak F_r$ is a polynomial of degree at most $r+1$ on the
affine hyperplane $\sum\alpha_j=1$. Its component of degree $r+1$
is entirely local and vanishes when $r>D$.
\end{theorem}
\begin{proof}
Substitute $U=s$ in \eqref{sp:germ} and extract $s^r$.
Homogeneous extraction in $e^{-L_{\boldsymbol\alpha}}$ gives
\eqref{sp:Rhomogeneous}. The degree assertions follow immediately.
\end{proof}

For a function on the normalized hyperplane, put
$\Delta_{\boldsymbol v}f(\boldsymbol\alpha)
=f(\boldsymbol\alpha+\boldsymbol v)-f(\boldsymbol\alpha)$ when
$\sum_jv_j=0$. The notation $[n^{-1}]f(n)$ means the coefficient
of $n^{-1}$ in the Laurent expansion at large $n$; it is the negative
of the complex-analytic residue at infinity when that residue applies.

\begin{theorem}[Local polarization identity]\label{sp:thm:polarization}
Let $\boldsymbol v_1,\ldots,\boldsymbol v_{r+1}$ be tangent vectors
with $\sum_jv_{ij}=0$, and put
$V_i(n)=\sum_jv_{ij}\log(n+a_j)$. Then
\begin{align}\label{sp:finite-difference}
 &\Delta_{\boldsymbol v_1}\cdots\Delta_{\boldsymbol v_{r+1}}
       \mathfrak F_r(\boldsymbol\alpha)\notag\\
 &\quad=D_{\boldsymbol v_1}\cdots D_{\boldsymbol v_{r+1}}
       \mathfrak F_r(\boldsymbol\alpha)
       =(-1)^{r+1}r!\,[n^{-1}]P(n)\prod_{i=1}^{r+1}V_i(n).
\end{align}
All mixed finite differences of order $r+2$ vanish. If $\deg P=r$
and $p$ is its leading coefficient, the right side reduces to
\begin{equation}\label{sp:leading}
 (-1)^{r+1}r!p
             \prod_{i=1}^{r+1}\left(\sum_jv_{ij}a_j\right).
\end{equation}
\end{theorem}

\begin{proof}
The $r+1$ tangent differentiations kill $\cH_{P,[r]}$ and polarize
$L_{\boldsymbol\alpha}^{r+1}$, supplying a factor $(r+1)!$ in
\eqref{sp:Rhomogeneous}. Since $\sum_jv_{ij}=0$,
\[
 L_{\boldsymbol v_i}((n+b)^{-1})
   =\sum_jv_{ij}\log(n+a_j)=V_i(n)=O(n^{-1}).
\]
The coefficient $\sum_k p_k[z^{k+1}]\prod_iL_{\boldsymbol v_i}(z)$
is therefore $[n^{-1}]P(n)\prod_i V_i(n)$. Translation from $n+b$
to $n$ preserves this coefficient. For a polynomial of degree at most
$r+1$, its mixed difference of order $r+1$ equals its corresponding
constant mixed derivative; for example integrate that derivative over
the unit $(r+1)$-cube. This proves \eqref{sp:finite-difference}.
Finally $V_i(n)=n^{-1}\sum_jv_{ij}a_j+O(n^{-2})$ proves
\eqref{sp:leading}.
\end{proof}

Normalization cannot be omitted. For arbitrary $\boldsymbol w$ with
$W\ne0$, define $Z_{\boldsymbol w}(s)=\cZ_P(s\boldsymbol w)$.
Then
\begin{equation}\label{sp:scaling}
 Z_{\boldsymbol w}^{(r)}(0)
       =W^r\mathfrak F_r(\boldsymbol w/W).
\end{equation}
The degree theorem concerns factor exponents; it is not a polynomiality
assertion in the shifts $a_j$ themselves.

\subsection{A full Hurwitz--Stieltjes formula for each jet}

Put $A_{h,\ell}=[z^\ell]L_{\boldsymbol\alpha}(z)^h$, with
$A_{0,0}=1$ and $A_{0,\ell}=0$ for $\ell>0$.

\begin{corollary}[Normally convergent coefficient formula]\label{sp:cor:alljets}
For every $r\ge0$,
\begin{align}\label{sp:alljets}
 \mathfrak F_r(\boldsymbol\alpha)
 =\sum_{k=0}^Dp_k\Bigg\{&\zeta^{(r)}(-k,b)
 +\sum_{\substack{\ell\ge1\\\ell\ne k+1}}\sum_{h=1}^r
   (-1)^h\binom rh A_{h,\ell}\zeta^{(r-h)}(\ell-k,b)\notag\\
 &+\frac{(-1)^{r+1}}{r+1}A_{r+1,k+1}
   +(-1)^r\sum_{h=1}^r\binom rh A_{h,k+1}\gamma_{r-h}(b)
                                      \Bigg\}.
\end{align}
Empty sums vanish. The generalized Stieltjes constants occur linearly,
with indices at most $r-1$.
\end{corollary}
\begin{proof}
Multiply the expansion of $e^{-sL_{\boldsymbol\alpha}(z)}$ by the
Taylor expansion of each nonresonant Hurwitz factor in
\eqref{sp:continuation}. For $\ell=k+1$, multiply it by
$s^{-1}+\sum_q(-1)^q\gamma_q(b)s^q/q!$. The coefficient of $s^r$
receives a contribution from the exponential coefficient of order
$r+1$, which yields the third term in \eqref{sp:alljets}. Normal
convergence was proved in Theorem~\ref{sp:thm:germ}.
\end{proof}

This additional order-$r+1$ coefficient must be retained before a
spectral pole is canceled. Generalized harmonic numbers provide an
alternative finite description of all binomial coefficients required:
\begin{equation}\label{sp:harmonic}
 \frac{(x)_q}{q!}
 =\frac xq\exp\left(\sum_{j\ge1}
                   \frac{(-1)^{j+1}H_{q-1}^{(j)}}j x^j\right),
 \qquad q\ge1.
\end{equation}
This follows from $(x)_q/q!=(x/q)\prod_{j=1}^{q-1}(1+x/j)$.
Use it in \eqref{sp:Cpoch} and truncate at the desired spectral order.
Thus both the Stieltjes and generalized-harmonic content are explicit.

\subsection{First derivative: the classical weighted anomaly and Gamma factors}

Set
\begin{equation}\label{sp:single}
 \zeta_P(s,a)=\sum_{n\ge0}\frac{P(n)}{(n+a)^s}
       =\sum_{k=0}^D\frac{P^{(k)}(-a)}{k!}\zeta(s-k,a),
\end{equation}
where continuation is understood, and define
$\delta_{\boldsymbol w}=\sum_jw_j\zeta_P'(0,a_j)
-Z_{\boldsymbol w}'(0)$. The first-derivative instance of the calculus
recovers
\begin{equation}\label{sp:anomaly}
 \boxed{\delta_{\boldsymbol w}
   =\frac1{2W}\sum_{i<j}w_iw_j[n^{-1}]P(n)
                 \log^2\frac{n+a_i}{n+a_j}.}
\end{equation}
Indeed the nonlocal terms of $\mathfrak F_1$ are linear in the
normalized weights and cancel against their interpolation at the
vertices $\boldsymbol e_j$. The local difference is
\[
 \frac12\sum_kp_k[z^{k+1}]
 \left\{\sum_j\alpha_jL_j^2-
           \left(\sum_j\alpha_jL_j\right)^2\right\}
 =\frac12\sum_{i<j}\alpha_i\alpha_j
              \sum_kp_k[z^{k+1}](L_i-L_j)^2.
\]
Multiplication by $W$ gives \eqref{sp:anomaly}.

The sign convention is fixed by the determinant identity
\begin{equation}\label{sp:determinant}
 e^{-Z_{\boldsymbol w}'(0)}
   =e^{\delta_{\boldsymbol w}}
                     \prod_j e^{-w_j\zeta_P'(0,a_j)}.
\end{equation}
For $P(n)=\sum_{d=0}^Dq_dn^d$, the finite coefficient in
\eqref{sp:anomaly} is
\begin{equation}\label{sp:pairpolynomial}
 \sum_{d=1}^Dq_d(-1)^{d+1}
   \sum_{\ell=1}^d
   \frac{(a_i^\ell-a_j^\ell)
         (a_i^{d+1-\ell}-a_j^{d+1-\ell})}{\ell(d+1-\ell)}.
\end{equation}
For $P=1$ it is zero; for $P=n,n^2,n^3$ it is, respectively,
\begin{align*}
 &(a_i-a_j)^2,\qquad -(a_i-a_j)^2(a_i+a_j),\\
 &\frac{(a_i-a_j)^2(11a_i^2+14a_ia_j+11a_j^2)}{12}.
\end{align*}
For shifts $(1,2,3)$ and weights $(1,2,1)$, this gives
$\delta=1$ for $P=n$ and $\delta=-4$ for $P=n^2$.

With $G$ denoting the Barnes $G$ function normalized by $G(1)=1$,
$G(a+1)=\Gamma(a)G(a)$, the Hurwitz--Barnes identity gives, when $P(n)=n$,
\begin{equation}\label{sp:Barnes}
 \zeta_P'(0,a)=\zeta'(-1)-\log G(a+1)+\frac a2\log(2\pi).
\end{equation}
For a direct normalization check, differentiate both sides using
Hurwitz shift differentiation and the Barnes integral formula
\cite[Eq.~5.17.4]{dlmf}. Both derivatives are
$a-1/2-a\psi(a)$, and their values at $a=1$ agree. This proves the
displayed relation on $(0,\infty)$.
For two unit-weight factors the complete determinant is therefore
\begin{equation}\label{sp:Barnesdet}
 e^{-Z_{(1,1)}'(0)}
 =\frac{e^{(a-b)^2/4-2\zeta'(-1)}G(a+1)G(b+1)}
             {(2\pi)^{(a+b)/2}}.
\end{equation}
Here $a,b>0$ are the two shifts; $b$ in this last formula need not be
the expansion base used earlier. All logarithms of $G$ are real on
these positive arguments.

\subsection{A six-term identity for second derivatives}

\begin{theorem}[Quadratic-weight Vandermonde evaluation]
\label{sp:thm:sixterm}
Let $P$ be a monic quadratic polynomial and, for $a,b>0$, define
\begin{equation}\label{sp:Tdef}
 \mathfrak T_P(a,b)=\left.\frac{d^2}{ds^2}
     \sum_{n\ge0}P(n)\bigl((n+a)^2(n+b)\bigr)^{-s}
                                      \right|_{s=0},
\end{equation}
by meromorphic continuation of the initially convergent series. Then
for every $a,b,c>0$,
\begin{align}\label{sp:sixterm}
 &\mathfrak T_P(a,b)+\mathfrak T_P(b,c)+\mathfrak T_P(c,a)
 -\mathfrak T_P(b,a)-\mathfrak T_P(c,b)-\mathfrak T_P(a,c)\notag\\
 &\hspace{30mm}=\frac23(a-b)(b-c)(c-a).
\end{align}
For a quadratic of leading coefficient $p$, multiply the right side
by $p$. For degree at most one, the right side is zero.
\end{theorem}

\begin{proof}
Use the three shifts $(a,b,c)$, take
$\boldsymbol\alpha_0=(1,1,1)/3$, and take tangent vectors
\[
 \boldsymbol v_1=(1,-1,0)/3,\quad
 \boldsymbol v_2=(0,1,-1)/3,\quad
 \boldsymbol v_3=(-1,0,1)/3.
\]
Their sum is zero, so the empty and full subsets cancel in the
third finite difference of $\mathfrak F_2$. Equation~\eqref{sp:leading}
gives
\[
 \sum_{i=1}^3\{\mathfrak F_2(\boldsymbol\alpha_0+\boldsymbol v_i)
      -\mathfrak F_2(\boldsymbol\alpha_0-\boldsymbol v_i)\}
 =-\frac2{27}(a-b)(b-c)(c-a).
\]
Each argument is a permutation of $(2/3,1/3,0)$. By
\eqref{sp:scaling} its value is the corresponding $\mathfrak T_P/9$.
Reading the six arguments gives the negative of the left side of
\eqref{sp:sixterm}, divided by nine. This proves the formula and its
sign. The lower-degree assertions follow from the coefficient
$[n^{-1}]$ in \eqref{sp:finite-difference}.
\end{proof}

For $(a,b,c)=(1,2,3)$ the value is $4/3$, independent of the lower
two coefficients of $P$. This is a cancellation of all nonlocal
second-derivative constants; applying only a first-derivative anomaly
formula would not justify it.

\subsection{An alternating Vandermonde identity with any number of factors}

The six-term evaluation is the first nontrivial higher-derivative member
of a general alternating family. Define
\[
 N=\binom m2,\qquad
 \Delta(\boldsymbol x)=\prod_{1\le i<j\le m}(x_j-x_i),\qquad
 K_m=\prod_{k=0}^{m-1}k!,
\]
and let $(\sigma\boldsymbol w)_j=w_{\sigma(j)}$.

\begin{theorem}[General spectral alternant]\label{sp:thm:alternant}
For $m\ge2$, shifts $a_j>0$, and weights with $W=\sum_jw_j\ne0$,
\begin{align}\label{sp:alternant}
 &\sum_{\sigma\in S_m}\operatorname{sgn}(\sigma)
       Z_{\sigma\boldsymbol w}^{(N-1)}(0)\notag\\
 &\quad=\frac{(-1)^N(N-1)!}{WK_m}\Delta(\boldsymbol w)
     [n^{-1}]P(n)\prod_{i<j}\log\frac{n+a_j}{n+a_i}.
\end{align}
If $\deg P=N-1$ with leading coefficient $p$, this is the complete
algebraic evaluation
\begin{equation}\label{sp:alternant-leading}
 \boxed{\sum_{\sigma\in S_m}\operatorname{sgn}(\sigma)
       Z_{\sigma\boldsymbol w}^{(N-1)}(0)
 =\frac{(-1)^N(N-1)!p}{WK_m}
                 \Delta(\boldsymbol w)\Delta(\boldsymbol a).}
\end{equation}
For degree less than $N-1$, the alternating sum is zero.
\end{theorem}

\begin{proof}
Every alternating polynomial in $m$ variables is divisible by their
Vandermonde polynomial, of degree $N$. Therefore alternation kills
the degree-at-most-$N-1$ term $\cH_{P,[N-1]}$ in
\eqref{sp:homogeneous}. The elementary alternant identity is
\begin{equation}\label{sp:elementary-alternant}
 \sum_{\sigma\in S_m}\operatorname{sgn}(\sigma)
       \left(\sum_j\alpha_{\sigma(j)}y_j\right)^N
   =\frac{N!}{K_m}\Delta(\boldsymbol\alpha)\Delta(\boldsymbol y).
\end{equation}
To prove it, expand by the multinomial theorem. A monomial with
repeated exponents disappears under alternation. The only distinct
nonnegative exponent pattern of total degree $N$ is a permutation
of $0,1,\ldots,m-1$. The surviving determinant in the $\alpha_j$
and the corresponding determinant in the $y_j$ give the two
Vandermondes, with the positive orientation shown in
\eqref{sp:elementary-alternant} and multinomial factor $N!/K_m$.

Apply this identity to $y_j=L_j(z)$ in
\eqref{sp:Rhomogeneous}. The normalized alternating jet is
\[
 \frac{(-1)^N(N-1)!}{K_m}\Delta(\boldsymbol\alpha)
           \sum_kp_k[z^{k+1}]\prod_{i<j}(L_j(z)-L_i(z)).
\]
At $z=(n+b)^{-1}$, every difference is
$\log((n+a_j)/(n+a_i))$. The coefficient is therefore the one
in \eqref{sp:alternant}. Use \eqref{sp:scaling} and
$\Delta(\boldsymbol w/W)=W^{-N}\Delta(\boldsymbol w)$ to obtain
the unnormalized factor $1/W$. Finally, the logarithmic product
starts with $\Delta(\boldsymbol a)n^{-N}$, proving
\eqref{sp:alternant-leading} and the lower-degree vanishing.
\end{proof}

For $m=3$ and $\boldsymbol w=(2,1,0)$, equation
\eqref{sp:alternant-leading} is exactly \eqref{sp:sixterm}.
For $m=4$, $N=6$ and $(N-1)!/K_m=10$: a signed sum of
twenty-four fifth spectral derivatives equals
$10p\Delta(\boldsymbol w)\Delta(\boldsymbol a)/W$.
For example, with $\boldsymbol w=(3,2,1,0)$,
$\boldsymbol a=(1,2,3,4)$ and any monic degree-five $P$, that
sum is exactly $240$. No individual fifth derivative is evaluated
to obtain this algebraic cancellation.

\subsection{Shift derivatives, polygammas, and normalized primitives}

The single-factor spectral derivative has the exact shift derivative
\begin{equation}\label{sp:shiftderivative}
 \frac{d}{da}\zeta_P'(0,a)
  =-\FP_{s=1}\zeta_P(s,a)
  =P(-a)\psi(a)
      -\sum_{k=1}^D\frac{P^{(k)}(-a)}{k!}\zeta(1-k,a).
\end{equation}
To prove it, differentiate the convergent Dirichlet series to obtain
$\partial_a\zeta_P(s,a)=-s\zeta_P(s+1,a)$. The latter zeta factor
has residue $P(-a)$ at $s=0$. Multiplication by $s$ before
differentiation shows that its finite part, rather than an undefined
value at the pole, is required. Since
$\zeta(1-k,a)=-B_k(a)/k$, the result is a polynomial times digamma
plus a polynomial. Repeated differentiation gives explicit polygamma
identities; integration over $[a_0,a_1]\subset(0,\infty)$ gives the
normalized primitive $\zeta_P'(0,a_1)-\zeta_P'(0,a_0)$.

Combining \eqref{sp:shiftderivative} with the finite anomaly polynomial
gives all shift derivatives of the product determinant. For example,
for $P=n$ and two unit weights,
\[
 \partial_a Z_{(1,1)}'(0)=\frac{a+b-1}2-a\psi(a).
\]

\subsection{Logarithmic multiplicities and higher-order poles}

Fix the same expansion base $b$ as in \eqref{sp:base}, and for an
integer $h\ge0$ consider the distinct spectrum
\begin{equation}\label{sp:logdef}
 \cZ_{P,h}(\boldsymbol u;\boldsymbol a;b)
  =\sum_{n\ge0}P(n)\log^h(n+b)\prod_j(n+a_j)^{-u_j}.
\end{equation}
Changing this $b$ changes the logarithmic multiplicity itself.
It must not be treated as a freely variable auxiliary base in this
definition.

\begin{theorem}[Higher-pole finite parts and their local identities]
\label{sp:thm:log}
Near zero,
\begin{equation}\label{sp:loggerm}
 \cZ_{P,h}(\boldsymbol u;\boldsymbol a;b)
   =\frac{h!\cR_P(\boldsymbol u)}{U^{h+1}}
                                +\cH_{P,h}(\boldsymbol u),
\end{equation}
where $\cH_{P,h}$ is holomorphic. For normalized $\boldsymbol\alpha$,
define the Laurent jet
\begin{equation}\label{sp:logjet}
 \mathfrak J_{r,h}(\boldsymbol\alpha)
  :=r![s^r]\cZ_{P,h}(s\boldsymbol\alpha;\boldsymbol a;b),
 \qquad r\ge0.
\end{equation}
It is a polynomial of degree at most $q:=r+h+1$, and
\begin{align}
 \mathfrak J_{r,h}
   &=r!\{\cH_{P,h,[r]}+h!\cR_{P,[r+h+1]}\},\label{sp:loghomogeneous}\\
 \Delta_{\boldsymbol v_1}\cdots\Delta_{\boldsymbol v_q}
       \mathfrak J_{r,h}(\boldsymbol\alpha)
   &=h!(-1)^q r!\,[n^{-1}]P(n)\prod_{i=1}^q V_i(n)
 \label{sp:logdifference}
\end{align}
for tangent vectors as before. All differences of order $q+1$ vanish.
\end{theorem}

\begin{proof}
The proof of \eqref{sp:continuation} is unchanged, except that each
Hurwitz factor is replaced by $(-1)^h\zeta^{(h)}$. At resonance,
\[
 (-1)^h\zeta^{(h)}(1+U,b)
  =\frac{h!}{U^{h+1}}
          +\sum_{j\ge0}\frac{(-1)^j\gamma_{j+h}(b)}{j!}U^j.
\]
Subtract the displayed pole to prove \eqref{sp:loggerm}. Normal
convergence persists after these fixed derivatives. Taking $[s^r]$
gives \eqref{sp:loghomogeneous}, and polarizing its degree-$q$ part
as in Theorem~\ref{sp:thm:polarization} proves
\eqref{sp:logdifference}.
\end{proof}

If $r+h>D$, the residue contribution in \eqref{sp:loghomogeneous}
vanishes, sharpening the degree bound to $r$. In particular, if $h>D$,
the constant finite part is independent of the factor shifts and
normalized weights and equals
$(-1)^h\sum_kp_k\zeta^{(h)}(-k,b)$, even though negative Laurent
powers may remain.

There is also a higher-pole version of Theorem~\ref{sp:thm:alternant}.
Let $N=\binom m2$, $0\le h\le N-1$, and $r=N-h-1$. The identical
alternant proof applied to \eqref{sp:loghomogeneous} gives
\begin{align}\label{sp:logalternant}
 &\sum_{\sigma\in S_m}\operatorname{sgn}(\sigma)\,
      r![s^r]\cZ_{P,h}(s\sigma\boldsymbol w;\boldsymbol a;b)\notag\\
 &\quad=\frac{h!(-1)^N r!}{W^{h+1}K_m}\Delta(\boldsymbol w)
       [n^{-1}]P(n)\prod_{i<j}\log\frac{n+a_j}{n+a_i}.
\end{align}
For $\deg P=N-1$ the final coefficient is
$p\Delta(\boldsymbol a)$. Thus the alternating algebraic evaluation
persists for logarithmic multiplicities, with the pole order determining
the required Laurent jet and the power of $W$.

For $h>0$ the ray generally has a pole of order at most $h$.
Consequently \eqref{sp:logjet} is a Laurent coefficient, not an
ordinary derivative of the unmodified spectrum at zero. With
$\mu=\sum_j\alpha_j(a_j-b)$, two concrete constant terms are
\begin{align}
 \FP_{s=0}\sum_{n\ge0}n\log(n+b)\prod_j(n+a_j)^{-\alpha_js}
  &=-\zeta'(-1,b)+b\zeta'(0,b)+\frac{\mu^2}2,
 \label{sp:logexample1}\\
 \FP_{s=0}\sum_{n\ge0}n^2\log^2(n+b)\prod_j(n+a_j)^{-\alpha_js}
  &=\zeta''(-2,b)-2b\zeta''(-1,b)+b^2\zeta''(0,b)
       -\frac{\mu^3}3.
 \label{sp:logexample2}
\end{align}
The positive-degree dependence of these finite parts is completely
algebraic despite the higher spectral poles. These formulas resolve
the concrete logarithmic-multiplicity extension of the local calculus;
arbitrary nonpolynomial spectra require a separate continuation theorem.


% END sections/07-spectral-products.tex
% BEGIN sections/08-audit-research.tex
\section{Audit, integration status, and further research}
\label{sec:audit}

\subsection{Precise disposition of the source questions}

The source baseline matters because the manuscript has already absorbed
many earlier conjectures and corrections. The audit below concerns the
relevant canonical sections and the five incoming reports at the pinned
revision; it is not a claim that every line of the collective manuscript
has been independently re-proved.

\begin{longtable}{@{}>{\raggedright\arraybackslash}p{.25\textwidth}>{\raggedright\arraybackslash}p{.43\textwidth}>{\raggedright\arraybackslash}p{.26\textwidth}@{}}
\toprule Source target & Result of this continuation & Remaining scope\\\midrule
\endhead
Nested Harmonic Jets, question 7 &
Theorem~\ref{mz:thm:zero} constructs every $F_R(-a-m;a)$.
Theorems~\ref{mz:thm:primitive} and \ref{mz:thm:coefficients} make the
coordinate construction finite and explicit. &
The smallest arithmetic or differential-algebraic coordinate class is
not proved.\\[5pt]
Nested Harmonic Jets, question 4 &
Theorems~\ref{thm:doublegerm}--\ref{thm:curved} give every transverse
ordered depth-two finite part, including curvature and reparametrization.
Theorem~\ref{thm:symcoeff} gives every fully symmetrized direction. &
Individual ordered germs in higher depth and compatible renormalization
remain separate problems.\\[5pt]
Nested Harmonic Jets, question 3 &
Independent perturbation directions have a complete finite coefficient
algorithm after full labelled symmetrization. &
This does not construct the general two-colour ordered product or every
non-diagonal Bell model.\\[5pt]
Canonical Gaussian $S_6$ and current $S_8$ candidates &
Neither short arithmetic evaluation is proved here. The current
source's conjectural labels must remain. &
A relation outside the restricted fixed-weight product presentation is
needed.\\[5pt]
Inverse-colour parity &
Section~\ref{sec:gap} supplies a two-shift disk formula, its derivative
calculus, and complete complementary-diagonal reductions. &
The underlying general parity and cyclotomic symmetry principles have
prior proofs.\\[5pt]
Polynomial spectral products &
Section~\ref{sec:spectral} proves an all-order normalized-weight calculus
and local identities, including higher poles. &
The first-derivative multiplicative anomaly retains its classical
attribution.\\
\bottomrule
\end{longtable}

\paragraph{Why the Gaussian candidates have not been promoted.}
The canonical manuscript contains an exact obstruction in its specified
level-four product-row presentation: a separating functional vanishes on
the Gaussian-supported rows but not on the mixed $S_{2m}$ coordinate
\cite[\src{chapters/05-level4-rank.tex}]{repo}.
The already proved $S_4$ identity uses a larger vocabulary, so this is
a proof-system obstruction, not an arithmetic independence theorem.
The present shifted-gap reflection belongs to the established parity
mechanism; simply specializing it does not supply the missing
weight-seven or weight-nine certificate. The source also distinguishes
its current frozen $S_8$ vector from a previously rejected vector.
Neither a tiny residual interval containing zero nor equivalence of two
candidate baskets proves the vanishing of the residual.

\subsection{A confirmed external indexing correction}

There is one concrete typographical correction in the external material.
In G\"urel \cite[Theorem 1.1, printed p.~2]{gurel}, arXiv version 2 of
3 September 2025, the second exponential sum is indexed by $\ell$ but
its denominator is printed as $k$. The original PDF, not only the HTML
conversion, contains this index. The denominator should be $\ell$.

Here is a direct derivation that fixes the correction. Under the
regularizable-spectrum hypotheses of that theorem, write
$R_\ell=\operatorname{Res}_{s=\ell}\zeta_\Lambda(s)$ and
$C_\ell=\FP_{s=\ell}\zeta_\Lambda(s)$. Near $s=0$,
\[
 \frac{(s)_\ell}{\ell!}
       =\frac{s}{\ell}+\frac{H_{\ell-1}s^2}{\ell}+O(s^3),
 \qquad
 \zeta_\Lambda(\ell+s)=\frac{R_\ell}{s}+C_\ell+O(s).
\]
The coefficient of $s$ in their product is
$(C_\ell+H_{\ell-1}R_\ell)/\ell$. Expanding the shifted spectral
series and negating its derivative therefore gives the corrected form
\begin{equation}\label{eq:Gurel-correction}
 D_\Lambda(z)=\exp\left\{-\zeta_\Lambda'(0)
  -\sum_{\ell=1}^{M}
     \frac{C_\ell+H_{\ell-1}R_\ell}{\ell}z^\ell\right\}
                                                    \Delta_\Lambda(z),
\end{equation}
where $H_0=0$, and $M$ and the canonical product $\Delta_\Lambda$
have the source's meaning. The convergent terms with $\ell>M$ form
the logarithm of that canonical product. This is an indexing correction,
not a refutation of the theorem.

\subsection{Recommended clarifications for integration}

No fresh substantive error was found in the relevant current repository
proofs examined here. The following clarifications prevent errors when
the new sections are integrated.

\begin{enumerate}[leftmargin=*,itemsep=4pt]
\item \textbf{Specify the regulator before a finite part.}
Replace any informal reference to ``the finite part at $(1,1)$'' by an
explicit path or subtraction convention. Equations
\eqref{eq:direction-FP} and \eqref{eq:curved-FP} are concrete transition
formulas between these choices.
\item \textbf{Retain the labelled symmetry factor.}
The fully symmetrized depth-$d$ value has $d!$ copies on the diagonal.
The source's individual diagonal convention is recovered by dividing by
$d!$, not by suppressing repeated permutations.
\item \textbf{Distinguish fixed and moving Gamma arguments.}
The value $F_R(-a-m;a)$ holds $z$ fixed while differentiating $s$.
The path $z=-(a+m)^s$ gives zero identically. Their derivative constraints
are compatible but are different statements.
\item \textbf{Multiply through a pole before differentiating.}
The spectral jet of order $r$ can require a coefficient of order $r+1$;
with pole order $h+1$ it can require order $r+h+1$. Equations
\eqref{sp:alljets} and \eqref{sp:loghomogeneous} show the missing terms
that a premature truncation would discard.
\item \textbf{Keep ordinary sums grouped and primitives normalized.}
The terms of \eqref{eq:product-summand} are summed together.
The combined right side of \eqref{mz:Mrec} is integrated together.
The index-one primitive \eqref{gap:eq:log-primitive} subtracts a base
point before summation. These are convergence statements in the
theorems, not discretionary numerical conventions.
\item \textbf{Preserve scope and attribution.}
Do not relabel general parity, the symmetric-sum identity, or the first
multiplicative anomaly as newly proved conjectures. Do not turn the
finite coordinate ceiling into a minimality theorem. Keep the current
$S_6$ and $S_8$ status unchanged unless a separate exact certificate is
supplied.
\end{enumerate}

\subsection{Reproducible checks and their evidential role}

The accompanying verification programs implement independent
representations where practical. The Gamma-zero calculations compare
the coordinate formula with Cauchy extraction from the original
spectral product. The shifted-gap checks compare bilateral reductions
with a Mellin integral. The anomaly checks use a finite initial segment
and an independently expanded Hurwitz tail before Cauchy extraction.
Exact polynomial checks test the finite combinatorial identities and
normalizations. The complete execution parameters, counts, residuals,
and package versions are retained in the machine-readable results.

The main numerical checks include sixteen Gamma-zero evaluations
through derivative order four, seven primitive-recursion checks,
four general shifted-gap reflections and three half-diagonal cases,
four first-derivative anomaly evaluations, and a six-term second-jet
evaluation. Observed absolute residuals in these independently evaluated
special-function checks are below $10^{-39}$; the recorded results give
the individual values and precision. These are floating-point
diagnostics, not rigorous interval certificates. Exact symbolic checks
and the analytic proofs have separate roles.

The package contains modular \LaTeX{} sources, the compiled article,
verification code, recorded results, a source-provenance manifest,
and integration notes with theorem labels. It uses no external files
from the source repository to compile. The original source snapshot is
identified by hashes rather than copied wholesale into the deliverable.

\subsection{Further research questions with concrete first targets}
\label{sec:questions}

\begin{enumerate}[leftmargin=*,label=\textbf{Q\arabic*.},itemsep=8pt]
\item \textbf{Ordered depth three, with independent exponents.}
Construct an analogue of \eqref{eq:doublegerm} for
$Z_3(1+u,1+v,1+w;a)$, displaying the intersecting polar divisors and
a normally convergent holomorphic remainder. A first deliverable is
an explicit constant term along every transverse ray, with its
symmetrization checked against \eqref{eq:sym3}. The comparison should
use the existing Laurent-expansion methods of \cite{mow}.

\item \textbf{A compatible subtraction algebra.}
Determine the local counterterms that convert the directional constants
of Sections~\ref{sec:direction}--\ref{sec:symmetric} into a specified
stuffle-preserving renormalization. The depth-two discrepancy
$-\gamma_1(a)(c/d+d/c)$ is the first necessary compatibility test.
Compare the result directly with \cite{guozhang}; path independence
must be a proved property of the chosen prescription.

\item \textbf{Tangential and nonlinear approaches to nested poles.}
Theorem~\ref{thm:curved} handles transverse analytic curves. Classify
curves tangent to $u=0$ or $u+v=0$ by their vanishing orders and give
the exact number of curve coefficients required for a finite part.
Curves lying identically in a polar divisor need an independently
specified restriction or subtraction operation.

\item \textbf{Smaller arithmetic classes for Gamma-zero coordinates.}
For a specified positive rational $x$, decide whether $D(x)$ and the
first higher coordinates $L_2(x),M_{2,2}(x)$ reduce to a stated field
of Gamma, polygamma, zeta-derivative, Dirichlet-$L$-derivative, and
algebraic-argument polylogarithm values. Fix the comparison field
before searching. Proving irreducibility requires more than failure
of symbolic simplification or an integer-relation search.

\item \textbf{Several isolated zero factors.}
Develop the counterpart of Theorem~\ref{mz:thm:zero} for a
multicolour product in which two or more factors vanish at the same
parameter. Determine how the multiplicity shifts the first nonzero
normal derivative, and how mixed spectral derivatives are encoded
by multivariate Bell polynomials. A two-colour periodic assignment
is the first case connecting directly to the source's question 3.

\item \textbf{The nonsymmetric Stieltjes sums.}
Equation~\eqref{eq:Stwo} evaluates the symmetric combination
$S_{m,n}+S_{n,m}$. Determine useful exact formulas for the difference
$S_{m,n}-S_{n,m}$, beginning with $(m,n)=(0,1)$. Through
\eqref{eq:Hcoeff}, this is exactly a question about the nonsymmetric
regular double-Hurwitz jet. Identify which part reduces to ordinary
Stieltjes constants and which part needs a new normalized integral.

\item \textbf{Primitive bases for complementary-shift polylogarithms.}
For $a\ge2$, integrate \eqref{gap:eq:diagonal} against $dz/z$ in a
specified classical or multiple-polylogarithm basis. The half-shift
formula \eqref{gap:eq:primitive} fixes the first normalization.
At rational shifts, track the root-of-unity alphabet explicitly and
separate variable-argument identities from isolated constant reductions.

\item \textbf{Noninteger indices and global shift continuation.}
Construct a contour or Mellin continuation of the two-shift reflection
in noninteger exponents. An integer-index partial fraction with a
finite summation range cannot be differentiated spectrally without
such an extension. Separately continue the shifts beyond the safe
strip, recording finite terms when lattice poles cross a contour.

\Needspace{5\baselineskip}
\item \textbf{Spectral interpolation and representation components.}
Classify finite identities between order-$r$ jets sampled at rational
normalized weights. Theorem~\ref{sp:thm:polynomial} turns this into
finite polynomial interpolation with a known top-degree correction.
Beyond alternating sums, determine which permutation-representation
components kill every nonlocal term and yield factorizations in
Vandermonde or Schur polynomials.

\item \textbf{Arithmetic and nonpolynomial multiplicities.}
Extend the all-order local calculus to quasi-polynomial or periodic
weights by residue-class splitting, and then to spectra whose base
zeta has several or higher-order poles. Section~\ref{sec:spectral}
already settles the concrete polynomial times logarithmic class.
For new classes, state the continuation and normal-convergence
hypotheses that justify every extracted coefficient.

\item \textbf{An exact certificate for $S_6$, then the current $S_8$.}
Construct a weight-seven functional equation using a larger Cayley or
iterated-integral vocabulary, reduce it to the frozen $S_6$ basket,
and exhibit a finite symbolic certificate. The source's separating
functional is a useful test that a genuinely additional relation has
entered the proof. Only after that mechanism is understood should
the weight-nine candidate be attacked by the same method. More
digits alone do not resolve the structural obstruction.

\item \textbf{Formal interfaces for finite identities and analytic limits.}
Formalize the subset telescoping law, partition M\"obius inversion,
Stirling rational-polylogarithm identity, Bell extraction, and tangent
polarization as exact algebraic components. Then formalize normal
convergence and Euler--Maclaurin remainder arguments as separate
analytic lemmas. Interval implementations of the independent numerical
checks would provide certified regression tests without being confused
with the symbolic proofs.
\end{enumerate}


% END sections/08-audit-research.tex
% BEGIN references.tex
\begin{thebibliography}{99}
\addcontentsline{toc}{section}{References}

\bibitem{repo}
V. Reshetnikov and contributors,
\emph{ProveIt: Polylogarithms and their Arithmetic Bridges}, canonical
manuscript and incoming research archives, repository revision
\src{16c7e342d7a4a15f59911f322eb90b7f96c2a8b5}, inspected
10 October 2026.
\url{https://github.com/VladimirReshetnikov/ProveIt/tree/16c7e342d7a4a15f59911f322eb90b7f96c2a8b5/Analysis/Polylogarithms/docs/manuscript}.

\bibitem{nested}
ProveIt incoming research continuation,
\emph{Nested Harmonic Jets and Gamma Renormalization},
10 October 2026, archive
\src{ProveIt_Nested_Harmonic_Jets_2026-10-10.zip} in \src{docs/incoming}
of \cite{repo}. Especially the zero-flow proposition and further
research questions 3, 4, and 7.

\bibitem{gauss}
ProveIt incoming research continuation,
\emph{Gauss--Hurwitz Pole Cancellation: Resonant Gamma-Ratio Sums,
Harmonic Identities, Stieltjes Jets, and Polylogarithmic Endpoint
Subtraction}, 10 October 2026, archive
\src{ProveIt_Gauss_Hurwitz_2026-10-10.zip} in \cite{repo}.

\bibitem{mellin}
ProveIt incoming research continuation,
\emph{Mellin--Lerch Transforms and Dilation Identities for Stieltjes
Correlations}, 10 October 2026, archive
\src{ProveIt_Mellin_Dilation_2026-10-10.zip} in \cite{repo}.

\bibitem{stharm}
ProveIt incoming research continuation,
\emph{Exact Stieltjes and Harmonic Identities}, 10 October 2026,
archive \src{ProveIt_Stieltjes_Harmonic_Identities_2026-10-10.zip}
in \cite{repo}.

\bibitem{translation}
ProveIt incoming research continuation,
\emph{Further Identities and Critical Transitions in the
Polylogarithm--Stieltjes Calculus}, 10 October 2026, archive
\src{ProveIt_Polylogarithm_Stieltjes_Continuation_2026-10-10.zip}
in \cite{repo}.

\bibitem{dlmf}
NIST, \emph{Digital Library of Mathematical Functions},
Sections 5.8 (Gamma products), 5.11 (asymptotics), 5.17 (Barnes $G$),
25.11 (Hurwitz zeta), and 25.12 (polylogarithms).
\url{https://dlmf.nist.gov/}. Accessed 10 October 2026.

\bibitem{hoffman}
M. E. Hoffman, Multiple harmonic series,
\emph{Pacific Journal of Mathematics} \textbf{152} (1992), 275--290.
\url{https://doi.org/10.2140/pjm.1992.152.275}.

\bibitem{mow}
K. Matsumoto, T. Onozuka, and I. Wakabayashi,
\emph{Laurent series expansions of multiple zeta-functions of
Euler--Zagier type at integer points}, 2016.
\url{https://arxiv.org/abs/1601.05918}.

\bibitem{guozhang}
L. Guo and B. Zhang, Renormalization of multiple zeta values,
\emph{Journal of Algebra} \textbf{319} (2008), 3770--3809.
\url{https://arxiv.org/abs/math/0606076}.

\bibitem{panzer}
E. Panzer, The parity theorem for multiple polylogarithms,
\emph{Journal of Number Theory} \textbf{172} (2017), 93--113.
\url{https://doi.org/10.1016/j.jnt.2016.08.004};
\url{https://arxiv.org/abs/1512.04482}.

\bibitem{rui}
H. Rui, \emph{Contour Integrations and Parity Results of Hurwitz-type
Cyclotomic Euler Sums}, preprint, submitted 30 December 2025.
\url{https://arxiv.org/abs/2601.00035}.

\bibitem{xu}
C. Xu, \emph{Symmetry Results for Cyclotomic Multiple Hurwitz Zeta
Values via Contour Integrals}, preprint, 2026; version 2,
10 September 2026.
\url{https://arxiv.org/abs/2602.10391}.

\bibitem{dowker}
J. S. Dowker, \emph{Calculation of the multiplicative anomaly},
2014, corrected version 3, 24 September 2023.
\url{https://arxiv.org/abs/1412.0549}.

\bibitem{gurel}
E. G\"urel, \emph{Resolution Of Multiplicative Anomaly Of Zeta
Regularization For Polynomials}, version 2, 3 September 2025.
\url{https://arxiv.org/pdf/2504.14563v2}.
The indexing correction in Theorem 1.1 is discussed in
Section~\ref{sec:audit} of the present article.

\bibitem{mizuno}
Y. Mizuno, Generalized Lerch formulas: Examples of zeta-regularized
products, \emph{Journal of Number Theory} \textbf{118} (2006), 155--171.
\url{https://doi.org/10.1016/j.jnt.2005.08.005}.

\end{thebibliography}

% END references.tex
\end{document}
